% ===================================================================
% LAPORAN PROYEK 1 - SHOWDRIVE
% Rancang Bangun Sistem Informasi Inventaris dan Transaksi
% Penjualan Kendaraan Bermotor Berbasis Web
%
% Program Studi D4 Teknik Informatika
% Universitas Logistik dan Bisnis Internasional (ULBI)
% Kelompok 18 — Tahun Akademik 2025/2026
%
% COMPILE: XeLaTeX (Overleaf → Menu → Compiler → XeLaTeX)
% ===================================================================

\documentclass[11pt, a5paper]{article}

% ---------------------------------------------------------------
% LAYOUT — A5, margin sesuai panduan ULBI
% ---------------------------------------------------------------
\usepackage[a5paper, top=2.54cm, bottom=2.54cm,
            left=3.17cm, right=3.17cm]{geometry}

% ---------------------------------------------------------------
% FONT (XeLaTeX — tersedia di semua distribusi Overleaf)
% ---------------------------------------------------------------
\usepackage{fontspec}
\setmainfont{Latin Modern Roman}
\setsansfont{Latin Modern Sans}
\setmonofont[Scale=0.88]{Latin Modern Mono}

% ---------------------------------------------------------------
% BAHASA INDONESIA
% ---------------------------------------------------------------
\usepackage[bahasai]{babel}

% ---------------------------------------------------------------
% FORMATTING PACKAGES
% ---------------------------------------------------------------
\usepackage{setspace}
\usepackage{titlesec}
\usepackage{tocloft}
\usepackage{tabularx}
\usepackage{booktabs}
\usepackage{multirow}
\usepackage{array}
\usepackage{graphicx}
\usepackage{float}
\usepackage{fancyhdr}
\usepackage{enumitem}
\usepackage{xcolor}
\usepackage{mdframed}
\usepackage{chngcntr}
\usepackage{amsmath}
\usepackage{amssymb}
% hyperref harus dimuat terakhir di antara packages non-TikZ
\usepackage[hidelinks]{hyperref}

% ---------------------------------------------------------------
% TIKZ
% ---------------------------------------------------------------
\usepackage{tikz}
\usepackage{dirtree}
\usetikzlibrary{
  shapes.geometric,
  shapes.misc,
  arrows.meta,
  positioning,
  fit,
  calc,
  backgrounds
}

% \Tiny tidak ada di LaTeX standar — definisikan manual
\makeatletter
\newcommand{\Tiny}{\@setfontsize\Tiny{5}{6}}
\makeatother

% --- STYLE GLOBAL DIAGRAM ---
% Node kotak biasa
\tikzstyle{box} = [
  rectangle, rounded corners=3pt,
  minimum width=2.2cm, minimum height=0.7cm,
  text centered, draw=black, fill=white,
  font=\small
]
% Node kotak besar (layer/tier)
\tikzstyle{bigbox} = [
  rectangle, rounded corners=4pt,
  minimum width=3.4cm, minimum height=1cm,
  text centered, draw=black!70, fill=gray!10,
  font=\small\bfseries
]
% Node flowchart proses
\tikzstyle{proc} = [
  rectangle, rounded corners=4pt,
  minimum width=3.5cm, minimum height=0.75cm,
  text centered, draw=black!80, fill=blue!8,
  font=\footnotesize, text width=3.2cm
]
% Node flowchart keputusan
\tikzstyle{dec} = [
  diamond, aspect=2.2,
  minimum width=3.2cm, minimum height=0.8cm,
  text centered, draw=black!80, fill=yellow!15,
  font=\footnotesize, text width=2.6cm
]
% Node start/end (oval)
\tikzstyle{term} = [
  ellipse, minimum width=2.6cm, minimum height=0.7cm,
  text centered, draw=black!80, fill=black!10,
  font=\footnotesize\bfseries
]
% Panah standar
\tikzstyle{arr}  = [-{Stealth[length=5pt]}, thick]
% Panah putus
\tikzstyle{darr} = [-{Stealth[length=5pt]}, thick, dashed]
% Panah dua arah
\tikzstyle{barr} = [{Stealth[length=5pt]}-{Stealth[length=5pt]}, thick]
% Node class diagram
\tikzstyle{cls}  = [
  rectangle, draw=black!80, fill=blue!5,
  minimum width=2.8cm, font=\footnotesize,
  text width=2.6cm, align=center
]
% Node ERD entitas
\tikzstyle{ent}  = [
  rectangle, rounded corners=2pt,
  draw=black!80, fill=green!8,
  minimum width=2.6cm, minimum height=0.7cm,
  text centered, font=\footnotesize\bfseries
]
% Node wireframe elemen UI
\tikzstyle{ui}   = [
  rectangle, rounded corners=2pt,
  draw=gray!70, fill=gray!10,
  text centered, font=\scriptsize,
  minimum height=0.55cm
]
\tikzstyle{uibtn}= [
  rectangle, rounded corners=3pt,
  draw=black!60, fill=black!80, text=white,
  text centered, font=\scriptsize\bfseries,
  minimum height=0.5cm
]
\tikzstyle{uicard}=[
  rectangle, rounded corners=3pt,
  draw=gray!60, fill=white,
  font=\scriptsize,
  text width=1.5cm, align=center,
  minimum height=1.4cm
]

% ---------------------------------------------------------------
% CAPTION — format "Gambar N. Judul" sesuai DOCX referensi
% ---------------------------------------------------------------
\usepackage[labelsep=period, figurename=Gambar,
            tablename=Tabel, font=small, labelfont=bf]{caption}

% ---------------------------------------------------------------
% JARAK BARIS
% ---------------------------------------------------------------
\setstretch{1.15}

% ---------------------------------------------------------------
% NONAKTIFKAN PENOMORAN OTOMATIS (BAB X dikontrol manual)
% ---------------------------------------------------------------
\setcounter{secnumdepth}{0}

% ---------------------------------------------------------------
% FORMAT JUDUL BAB / SUB-BAB / SUB-SUB-BAB
% ---------------------------------------------------------------
\titleformat{\section}[block]
  {\normalfont\Large\bfseries\centering}{}{}{}
\titleformat{\subsection}[block]
  {\normalfont\large\bfseries}{}{}{}
\titleformat{\subsubsection}[block]
  {\normalfont\normalsize\bfseries}{}{}{}

\titlespacing*{\section}   {0pt}{2.8ex plus 0.5ex minus 0.2ex}{1.6ex}
\titlespacing*{\subsection}{0pt}{2.0ex plus 0.3ex minus 0.2ex}{1.0ex}
\titlespacing*{\subsubsection}{0pt}{1.4ex plus 0.2ex}{0.7ex}

% ---------------------------------------------------------------
% LIST STYLE
% ---------------------------------------------------------------
\setlist[itemize]{label=\textbullet, leftmargin=*, itemsep=2pt}
\setlist[enumerate]{leftmargin=*, itemsep=2pt}

% ---------------------------------------------------------------
% NOMOR GAMBAR GLOBAL (tidak reset tiap section)
% ---------------------------------------------------------------
\counterwithout{figure}{section}

% ---------------------------------------------------------------
% DAFTAR ISI & GAMBAR
% ---------------------------------------------------------------
\renewcommand{\contentsname}{DAFTAR ISI}
\renewcommand{\listfigurename}{DAFTAR GAMBAR}
\renewcommand{\cftsecleader}{\cftdotfill{\cftdotsep}}

% ---------------------------------------------------------------
% KOTAK INFO — untuk catatan teknis penting (box highlight)
% ---------------------------------------------------------------
\newmdenv[
  backgroundcolor=gray!8,
  linecolor=gray!40,
  linewidth=0.6pt,
  innerleftmargin=10pt,
  innerrightmargin=10pt,
  innertopmargin=8pt,
  innerbottommargin=8pt,
  skipabove=8pt,
  skipbelow=8pt
]{infobox}

% ---------------------------------------------------------------
% HEADER & FOOTER
% ---------------------------------------------------------------
\pagestyle{plain}

% ===================================================================
\begin{document}
% ===================================================================

% ================================================================
%  COVER / HALAMAN JUDUL
% ================================================================
\begin{titlepage}
  \centering
  \vspace*{0.4cm}

  % Logo placeholder — ganti dengan \includegraphics{images/logo-ulbi}
  \fbox{\parbox{2.8cm}{\centering\vspace{1.3cm}
    {\small LOGO ULBI}\vspace{1.3cm}}}

  \vspace{0.5cm}
  {\large\bfseries UNIVERSITAS LOGISTIK DAN BISNIS INTERNASIONAL\par}
  \vspace{0.1cm}
  {\normalsize PROGRAM STUDI D4 TEKNIK INFORMATIKA\par}
  \vspace{0.3cm}
  \rule{\textwidth}{1pt}
  \vspace{0.3cm}

  {\Large\bfseries LAPORAN PROYEK 1\par}
  \vspace{0.5cm}
  {\large\bfseries
    RANCANG BANGUN SISTEM INFORMASI INVENTARIS\\[4pt]
    DAN TRANSAKSI PENJUALAN BERBASIS WEB\par}
  \vspace{0.25cm}
  {\small\itshape
    ShowDrive: Sistem Informasi Manajemen Showroom Otomotif Premium\par}

  \vspace{1.5cm}
  {\small Disusun oleh:\par}
  \vspace{0.4cm}
  {\large\bfseries Kelompok 18\par}
  \vspace{0.6cm}

  \begin{tabular}{rl}
    Fathoni Abdul Jabbar   & 714250039 \\[4pt]
    Muhammad Aris Ashidiqi & 714250042 \\
  \end{tabular}

  \vfill
  \rule{\textwidth}{0.4pt}
  \vspace{0.3cm}
  {\normalsize\bfseries BANDUNG, 2026\par}
\end{titlepage}

% ================================================================
%  BAGIAN AWAL — penomoran romawi, mulai halaman ii (cover = i)
% ================================================================
\pagenumbering{roman}
\setcounter{page}{2}

% ================================================================
%  KATA PENGANTAR
% ================================================================
\newpage
\section*{KATA PENGANTAR}
\addcontentsline{toc}{section}{KATA PENGANTAR}

Puji dan syukur penulis panjatkan ke hadirat Tuhan Yang Maha Esa atas rahmat
dan karunia-Nya sehingga Laporan Akhir Proyek 1 dengan judul \textit{``Rancang
Bangun Sistem Informasi Inventaris dan Transaksi Penjualan Berbasis Web''}
(ShowDrive) ini dapat diselesaikan dengan baik dan tepat waktu.

Laporan ini disusun sebagai pemenuhan syarat kelulusan mata kuliah Proyek 1
Program Studi D4 Teknik Informatika, Universitas Logistik dan Bisnis
Internasional (ULBI). Di samping itu, laporan ini dirancang sekaligus sebagai
\textbf{buku panduan teknis} yang komprehensif, sehingga pembaca — baik
mahasiswa, pengembang, maupun pelaku usaha showroom — dapat memahami cara kerja
sistem, filosofi desain, dan cara mengoperasikan aplikasi ShowDrive secara
mandiri.

Selama proses pengerjaan, penulis mengimplementasikan ilmu rekayasa perangkat
lunak, pemrograman web berbasis \textit{framework} Laravel, serta prinsip
perancangan basis data relasional untuk membangun sebuah platform manajemen
showroom yang aman, efisien, dan dapat dipercaya.

Penulis menyampaikan terima kasih yang sebesar-besarnya kepada:
\begin{enumerate}
  \item Dosen Pembimbing Kelompok 18 yang telah memberikan arahan, waktu, dan
        ilmu selama proses perancangan hingga implementasi sistem.
  \item Koordinator Mata Kuliah Proyek 1 yang telah memfasilitasi jalannya
        perkuliahan proyek secara terstruktur.
  \item Seluruh Dosen dan Staf Pengajar Program Studi D4 Teknik Informatika ULBI
        yang telah membekali pondasi keilmuan yang bermanfaat.
  \item Orang tua dan keluarga yang senantiasa memberikan dukungan moril dan
        materil.
  \item Seluruh rekan mahasiswa D4 Teknik Informatika angkatan 2025 yang saling
        mendukung selama pengerjaan proyek.
\end{enumerate}

Penulis menyadari laporan ini masih dapat dikembangkan lebih lanjut. Kritik dan
saran yang membangun sangat diharapkan demi kesempurnaan sistem di masa mendatang.

\vspace{1cm}
\begin{flushright}
  Bandung, Juli 2026\\[1.2cm]
  \textbf{Kelompok 18}
\end{flushright}

% ================================================================
%  LEMBAR PERSETUJUAN SIDANG
% ================================================================
\newpage
\begin{center}
  {\large\bfseries LEMBAR PERNYATAAN PERSETUJUAN}\\[4pt]
  {\large\bfseries DAN PERMOHONAN SIDANG PROYEK 1}
\end{center}
\addcontentsline{toc}{section}{LEMBAR PERSETUJUAN SIDANG}
\vspace{0.6cm}

Saya sebagai Pembimbing Kelompok 18 dengan Anggota:

\vspace{0.5cm}
\begin{center}
\renewcommand{\arraystretch}{1.7}
\begin{tabular}{|p{0.36\textwidth}|p{0.51\textwidth}|}
\hline
Nama Mahasiswa 1 & Fathoni Abdul Jabbar \\ \hline
NPM              & 714250039 \\ \hline
Nama Mahasiswa 2 & Muhammad Aris Ashidiqi \\ \hline
NPM              & 714250042 \\ \hline
Judul Proyek 1   & Rancang Bangun Sistem Informasi Inventaris dan
                   Transaksi Penjualan Berbasis Web \\ \hline
\end{tabular}
\end{center}
\vspace{0.5cm}

Menyatakan bahwa mahasiswa tersebut telah menyelesaikan semua luaran dengan
kemajuan: \dotfill\%. Bagian yang belum diselesaikan (Jika ada):

\vspace{0.15cm}
\noindent\dotfill

\vspace{0.1cm}
\noindent\dotfill

\vspace{0.4cm}
Adapun penulisan laporan Akhir Proyek 1 telah diselesaikan seluruhnya (100\%).
Dengan demikian saya mengajukan mahasiswa tersebut untuk mengikuti sidang Proyek
1. Apabila ternyata pernyataan saya tersebut tidak benar, maka saya menyetujui
penundaan sidang termasuk pembatalan sidang Proyek 1 untuk mahasiswa bimbingan
saya tersebut sesuai aturan yang berlaku.

\vspace{0.5cm}
Bandung, \hspace{3.5cm} 2026

\vspace{0.9cm}
\noindent
\begin{minipage}[t]{0.48\textwidth}
  \textbf{Mahasiswa 1}\\[2.2cm]
  \underline{Fathoni Abdul Jabbar}\\
  NPM: 714250039
\end{minipage}\hfill
\begin{minipage}[t]{0.48\textwidth}
  \textbf{Dosen Pembimbing}\\[2.2cm]
  \underline{\hspace{4.5cm}}\\
  NIK:\hspace{0.5cm}
\end{minipage}

\vspace{0.9cm}
\noindent
\begin{minipage}[t]{0.48\textwidth}
  \textbf{Mahasiswa 2}\\[2.2cm]
  \underline{Muhammad Aris Ashidiqi}\\
  NPM: 714250042
\end{minipage}

% ================================================================
%  LEMBAR PENGESAHAN
% ================================================================
\newpage
\begin{center}
  {\large\bfseries LEMBAR PENGESAHAN}\\[6pt]
  {\normalsize\bfseries RANCANG BANGUN SISTEM INFORMASI INVENTARIS}\\
  {\normalsize\bfseries DAN TRANSAKSI PENJUALAN BERBASIS WEB}
\end{center}
\addcontentsline{toc}{section}{LEMBAR PENGESAHAN}
\vspace{0.5cm}

\begin{flushleft}
  Fathoni Abdul Jabbar \hspace{1.5cm} \textit{714250039}\\[4pt]
  Muhammad Aris Ashidiqi \hspace{0.8cm} \textit{714250042}
\end{flushleft}

\vspace{0.4cm}
Dokumen Proyek 1 ini telah diperiksa, disetujui dan disidangkan\\
Di Bandung, \dotfill

\vspace{0.9cm}
Oleh:

\vspace{0.5cm}
\noindent
\begin{minipage}[t]{0.48\textwidth}
  Penguji Pendamping,\\[2.2cm]
  \underline{\hspace{4.5cm}}\\NIK:
\end{minipage}\hfill
\begin{minipage}[t]{0.48\textwidth}
  Penguji Utama,\\[2.2cm]
  \underline{\hspace{4.5cm}}\\NIK:
\end{minipage}

\vspace{0.9cm}
\noindent
\begin{minipage}[t]{0.48\textwidth}
  Pembimbing,\\[2.2cm]
  \underline{\hspace{4.5cm}}\\NIK:
\end{minipage}\hfill
\begin{minipage}[t]{0.48\textwidth}
  Koordinator Proyek 1,\\[2.2cm]
  \underline{\hspace{4.5cm}}\\NIK:
\end{minipage}

\vspace{0.9cm}
\noindent
Menyetujui,\\
Ketua Program Studi D-IV Teknik Informatika,\\[2.2cm]
\underline{\hspace{5cm}}\\NIK:

% ================================================================
%  ABSTRAK
% ================================================================
\newpage
\section*{ABSTRAK}
\addcontentsline{toc}{section}{ABSTRAK}

\noindent\textbf{Rancang Bangun Sistem Informasi Inventaris dan Transaksi
Penjualan Berbasis Web (ShowDrive)}\\[5pt]
\noindent\textit{Fathoni Abdul Jabbar (714250039) \& Muhammad Aris Ashidiqi
(714250042)}\\
\noindent\textit{Program Studi D4 Teknik Informatika, ULBI, 2026}

\vspace{0.5cm}

Pengelolaan inventaris kendaraan dan transaksi penjualan di showroom otomotif
konvensional masih banyak dilakukan secara manual, menggunakan lembar
\textit{spreadsheet} dan komunikasi tatap muka, sehingga rentan terhadap
ketidakakuratan data, duplikasi transaksi, dan rendahnya transparansi bagi
pelanggan. Penelitian ini merancang dan mengimplementasikan sistem informasi
manajemen showroom berbasis web bernama \textbf{ShowDrive} sebagai solusi
komprehensif atas permasalahan tersebut.

Sistem dibangun menggunakan \textit{framework} Laravel 12 (PHP 8.2) sebagai
\textit{backend} dengan basis data relasional MySQL 8.0 yang dinormalisasi ke
dalam 6 entitas utama. Fitur unggulan sistem mencakup: (1)~verifikasi identitas
pelanggan berbasis OTP 4 digit dengan sesi kedaluwarsa 5 menit dan batas 3
percobaan; (2)~dukungan dua skema pembayaran, yaitu \textit{Down Payment} (DP)
dengan persentase yang dikonfigurasi per unit dan pembayaran lunas; (3)~kalkulasi
Pajak Pertambahan Nilai (PPN) 11\% secara otomatis; (4)~mekanisme transaksi
atomik (\texttt{DB::transaction()}) yang menjamin konsistensi data lintas tabel;
serta (5)~integritas referensial berlapis dengan \textit{constraint ON DELETE}
berbeda untuk setiap relasi.

Hasil implementasi menunjukkan sistem berhasil memenuhi seluruh tujuan
pengembangan, melayani dua kelompok pengguna dalam satu platform terpadu dengan
jalur keamanan akses administrasi yang tersembunyi (\textit{security by
obscurity}).

\vspace{0.4cm}
\noindent\textbf{Kata Kunci:} Sistem Informasi Showroom, Laravel, Transaksi
Atomik, OTP, Down Payment, PPN, Integritas Referensial, \textit{Security by
Obscurity}.

% ================================================================
%  DAFTAR ISI
% ================================================================
\newpage
\tableofcontents

% ================================================================
%  DAFTAR GAMBAR
% ================================================================
\newpage
\listoffigures
\addcontentsline{toc}{section}{DAFTAR GAMBAR}

% ================================================================
%  PENOMORAN ARAB MULAI BAB ISI
% ================================================================
\newpage
\pagenumbering{arabic}
\setcounter{page}{1}

% ================================================================
%  BAB 1  PENDAHULUAN
% ================================================================
\section*{BAB 1 PENDAHULUAN}
\addcontentsline{toc}{section}{BAB 1 PENDAHULUAN}

% ----------------------------------------------------------------
\subsection*{1.1. Latar Belakang}
% ----------------------------------------------------------------

Industri otomotif merupakan salah satu sektor dengan volume transaksi tertinggi
dan nilai aset per transaksi yang sangat signifikan. Sebuah unit kendaraan
premium dapat bernilai miliaran rupiah, sehingga setiap kesalahan pencatatan
atau ketidakakuratan data inventaris berdampak langsung pada integritas keuangan
bisnis. Ironisnya, sebagian besar showroom kendaraan bermotor di Indonesia —
terutama skala menengah — masih mengandalkan pencatatan inventaris manual
menggunakan lembar kerja \textit{spreadsheet} atau bahkan catatan fisik, serta
proses pemesanan yang berlangsung sepenuhnya melalui komunikasi telepon atau
pertemuan langsung.

Setidaknya ada empat permasalahan utama yang secara konsisten ditemui dalam
model operasional konvensional tersebut. Pertama, \textbf{ketidakakuratan data
inventaris}: status ketersediaan unit kendaraan tidak selalu diperbarui secara
\textit{real-time}, sehingga pelanggan dapat melakukan pemesanan untuk unit yang
sebenarnya sudah terjual. Kedua, \textbf{minimnya transparansi transaksi}:
pelanggan tidak memiliki cara mandiri untuk memverifikasi apakah pembayaran yang
telah mereka lakukan sudah diterima dan diproses oleh pihak showroom, sehingga
harus menghubungi staf secara aktif. Ketiga, \textbf{proses verifikasi yang
lambat}: kasir harus memeriksa bukti transfer secara manual satu per satu dan
memperbarui status di lembar kerja terpisah, sebuah proses yang rawan kesalahan
dan memakan waktu saat volume transaksi tinggi. Keempat, \textbf{risiko duplikasi
dan kehilangan data}: tanpa sistem basis data yang ternormalisasi, riwayat
transaksi rentan terduplikasi atau hilang, terutama saat terjadi pergantian
perangkat atau staf yang mengelola data.

Menjawab keempat permasalahan tersebut, penelitian ini merancang dan
mengimplementasikan sistem informasi manajemen showroom kendaraan bermotor
berbasis web bernama \textbf{ShowDrive}. Sistem ini dibangun di atas
\textit{framework} Laravel 12 dengan PHP 8.2 sebagai \textit{backend}, dan
MySQL 8.0 sebagai basis data relasional yang dinormalisasi ke dalam enam entitas
utama. ShowDrive bukan sekadar aplikasi katalog, melainkan sebuah platform
transaksi terintegrasi yang mencakup seluruh siklus hidup pemesanan kendaraan
dari sisi pelanggan maupun sisi administrasi internal showroom.

\begin{infobox}
\textbf{Catatan Implementasi:} Sistem ShowDrive dirancang dan diuji pada
lingkungan \textit{server} lokal berbasis XAMPP dengan data sampel berupa unit
kendaraan premium (Porsche 911 GT3 RS, Chevrolet Corvette C8, dan Ferrari 488
Pista Coupe) guna mensimulasikan kondisi operasional showroom otomotif nyata.
\end{infobox}

% ----------------------------------------------------------------
\subsection*{1.2. Nama Aplikasi dan Dasar Ide}
% ----------------------------------------------------------------

Aplikasi yang dikembangkan dalam penelitian ini diberi nama \textbf{ShowDrive}.
Nama ini merupakan penggabungan dari dua kata dalam bahasa Inggris: \textit{Show}
(menampilkan atau memamerkan) dan \textit{Drive} (berkendara). Secara konseptual,
nama ini mencerminkan dua peran utama sistem: menampilkan katalog kendaraan secara
digital kepada publik, sekaligus memfasilitasi seluruh proses transaksi yang
berujung pada keputusan berkendara atau membeli kendaraan impian pelanggan.

Dasar ide pengembangan ShowDrive bermula dari pengamatan bahwa proses transaksi
showroom konvensional memiliki terlalu banyak titik kegagalan (\textit{failure
points}) yang seharusnya dapat dieliminasi dengan teknologi. Ide intinya adalah
memindahkan seluruh alur transaksi — dari penemuan unit kendaraan di katalog,
pengajuan pemesanan, pengunggahan bukti transfer, hingga penerbitan kwitansi
digital — ke dalam satu platform web yang dapat diakses kapan saja dan dari mana
saja, tanpa memerlukan kehadiran fisik pelanggan di showroom hingga tahap serah
terima unit.

Terdapat beberapa inovasi teknis yang menjadi ciri khas ShowDrive dan
membedakannya dari sistem manajemen showroom konvensional:

\begin{enumerate}[label=\alph*.]
  \item \textbf{Verifikasi Identitas Berbasis OTP:} Alih-alih sistem login
        berbasis akun dan kata sandi yang mengharuskan pelanggan mendaftar,
        ShowDrive menggunakan verifikasi OTP (\textit{One-Time Password}) 4 digit
        yang dikaitkan langsung dengan nomor HP pelanggan. OTP memiliki masa
        kedaluwarsa 5 menit dan dibatasi maksimal 3 kali percobaan verifikasi per
        sesi, kemudian di-\textit{hash} dengan Bcrypt sebelum disimpan di sesi
        untuk mencegah pembandingan nilai \textit{plaintext}.

  \item \textbf{Dukungan Skema Down Payment (DP):} Setiap unit kendaraan dapat
        dikonfigurasi dengan persentase DP tersendiri (misalnya 20\% atau 25\%).
        Pelanggan dapat memilih membayar penuh (\textit{Full Payment}) atau
        mencicil dengan uang muka terlebih dahulu. Sistem secara otomatis
        menghitung nominal DP berdasarkan harga unit dan persentase yang
        dikonfigurasi, sehingga tidak ada potensi kesalahan kalkulasi manual.

  \item \textbf{Kalkulasi PPN 11\% Otomatis:} Sesuai dengan regulasi perpajakan
        Indonesia, sistem secara otomatis menghitung dan mencatat Pajak
        Pertambahan Nilai (PPN) sebesar 11\% dari harga unit di setiap faktur
        yang diterbitkan. Nilai subtotal, nominal pajak, dan total akhir
        disimpan secara terpisah di basis data untuk kemudahan audit keuangan.

  \item \textbf{Perlindungan Integritas Data Tingkat Basis Data:} Sistem secara
        aktif mencegah penghapusan aset kendaraan yang masih terikat transaksi
        aktif menggunakan \textit{constraint ON DELETE RESTRICT} pada tingkat
        basis data MySQL, bukan hanya validasi di lapisan aplikasi.

  \item \textbf{Operasi Kasir Berbasis AJAX:} Seluruh aksi verifikasi, penolakan,
        pembatalan, amend jadwal, dan konfirmasi serah terima unit dilakukan
        secara asinkron tanpa \textit{page refresh}, dengan respons visual
        \textit{real-time} pada antarmuka kasir.
\end{enumerate}

% ----------------------------------------------------------------
\subsection*{1.3. Tujuan Pengembangan}
% ----------------------------------------------------------------

Pengembangan sistem ShowDrive bertujuan untuk menghasilkan sebuah aplikasi web
yang mampu:

\begin{enumerate}[label=\alph*.]
  \item \textbf{Mendigitalisasi Proses Pengelolaan Inventaris Kendaraan:}
        Menyediakan antarmuka yang memudahkan admin/kasir dalam menambahkan,
        memperbarui, dan menghapus data unit kendaraan beserta galeri foto
        secara terpusat, terstruktur, dan dilindungi dari penghapusan tidak sah
        melalui \textit{constraint} basis data.

  \item \textbf{Menyederhanakan Alur Pemesanan bagi Pelanggan:} Membangun
        pengalaman pemesanan yang intuitif dan sepenuhnya mandiri (\textit{self
        service}) — mulai dari penelusuran katalog, pengisian formulir pemesanan
        dengan validasi NIK opsional, kalkulasi harga otomatis dengan PPN,
        hingga pengunggahan bukti transfer — tanpa memerlukan registrasi akun
        terlebih dahulu.

  \item \textbf{Mengamankan Akses Pelacakan Transaksi:} Mengimplementasikan
        sistem OTP berbasis nomor HP untuk memverifikasi identitas pelanggan
        sebelum menampilkan riwayat transaksi, menggantikan sistem login
        konvensional yang mengharuskan pelanggan mengingat kata sandi.

  \item \textbf{Mempercepat Proses Validasi Transaksi oleh Kasir:}
        Mengimplementasikan fitur verifikasi pembayaran, penolakan dengan catatan,
        amend jadwal inspeksi, dan konfirmasi serah terima unit berbasis AJAX
        pada \textit{dashboard} administrasi sehingga kasir dapat memproses
        seluruh siklus transaksi dalam satu antarmuka terpadu secara
        \textit{real-time}.

  \item \textbf{Menjamin Keandalan dan Konsistensi Data Transaksi:} Menerapkan
        mekanisme transaksi basis data atomik (prinsip ACID) untuk memastikan
        bahwa operasi yang melibatkan beberapa tabel sekaligus berhasil
        sepenuhnya atau dibatalkan sepenuhnya, tanpa ada kondisi data yang
        setengah jadi.

  \item \textbf{Menyediakan Laporan Keuangan dan Metrik Operasional
        \textit{Real-Time}:} Menampilkan metrik performa showroom secara dinamis
        pada \textit{dashboard} admin, mencakup total unit inventaris, jumlah
        unit terjual, jumlah transaksi menunggu verifikasi, dan total pendapatan
        keseluruhan.
\end{enumerate}

% ----------------------------------------------------------------
\subsection*{1.4. Ruang Lingkup}
% ----------------------------------------------------------------

Untuk memastikan pengembangan sistem berjalan terarah, ruang lingkup proyek
ShowDrive dibatasi secara eksplisit sebagai berikut.

\subsubsection*{a. Aktor Pengguna}

Sistem ini dibatasi untuk 2 (dua) aktor utama dengan hak akses yang diatur
secara ketat:
\begin{itemize}
  \item \textbf{Pelanggan Publik} — mengakses website tanpa akun, diidentifikasi
        hanya melalui nomor HP yang diverifikasi dengan OTP.
  \item \textbf{Administrator/Kasir} — staf internal showroom yang masuk melalui
        jalur autentikasi aman dengan \textit{username} dan \textit{password}.
\end{itemize}

\subsubsection*{b. Fitur yang Tercakup}

\begin{itemize}
  \item Katalog kendaraan publik dengan fitur pencarian (\textit{full-text
        search} berdasarkan merek, model, dan VIN).
  \item Formulir pemesanan dengan validasi NIK opsional dan kalkulasi PPN 11\%
        otomatis.
  \item Dua skema pembayaran: \textit{Down Payment} (DP) dan pembayaran penuh.
  \item Verifikasi identitas pelanggan berbasis OTP sebelum akses riwayat
        transaksi.
  \item Pengunggahan bukti transfer (\textit{file} JPEG/PNG/WebP, maks.\ 2~MB;
        SVG diblokir untuk mencegah XSS).
  \item \textit{Dashboard} kasir dengan verifikasi pembayaran, penolakan, amend
        jadwal, dan konfirmasi serah terima unit via AJAX.
  \item Penerbitan invoice digital dengan format kode \texttt{SD/YYYY/NNNN}.
  \item Manajemen inventaris (CRUD kendaraan, galeri gambar, gudang, kasir,
        profil perusahaan, pengaturan rekening/QRIS).
  \item Laporan keuangan ringkas pada \textit{dashboard} admin.
\end{itemize}

\subsubsection*{c. Fitur di Luar Ruang Lingkup}

\begin{itemize}
  \item Integrasi \textit{payment gateway} otomatis (Midtrans, Xendit, dll.).
  \item Pengiriman notifikasi SMS atau WhatsApp secara otomatis melalui API
        pihak ketiga berbayar. (Fitur notifikasi WhatsApp disediakan sebagai
        \textit{prefilled message} yang dibuka secara manual oleh kasir.)
  \item Manajemen logistik pengiriman fisik kendaraan ke alamat pelanggan.
  \item Ekspor laporan keuangan ke format PDF atau Excel secara otomatis.
\end{itemize}

\subsubsection*{d. Platform dan Teknologi}

Sistem dibangun menggunakan Laravel 12 dengan PHP 8.2 (\textit{backend}), Blade
Template Engine dan Tailwind CSS (\textit{frontend}), serta MySQL 8.0 sebagai
basis data. Pengembangan dilakukan pada lingkungan server lokal XAMPP berbasis
Windows.

\newpage

% ================================================================
%  BAB 2  DESKRIPSI SISTEM
% ================================================================
\section*{BAB 2 DESKRIPSI SISTEM}
\addcontentsline{toc}{section}{BAB 2 DESKRIPSI SISTEM}

% ----------------------------------------------------------------
\subsection*{2.1. Gambaran Umum Aplikasi}
% ----------------------------------------------------------------

ShowDrive adalah platform manajemen showroom otomotif berbasis web yang
mengintegrasikan seluruh proses bisnis penjualan kendaraan — mulai dari
penayangan katalog unit, pemesanan oleh pelanggan, pengelolaan pembayaran,
hingga serah terima fisik kendaraan — ke dalam satu sistem informasi terpusat.
Sistem ini beroperasi pada basis data MySQL 8.0 dengan enam tabel relasional yang
saling terhubung, menjamin konsistensi dan integritas data di seluruh siklus
transaksi.

Secara arsitektur, ShowDrive melayani dua kelompok pengguna dengan antarmuka yang
terpisah namun berbagi basis data yang sama:

\begin{enumerate}
  \item \textbf{Antarmuka Publik (Pelanggan):} Dapat diakses oleh siapa saja
        tanpa perlu mendaftar akun. Pelanggan menelusuri katalog kendaraan,
        memilih unit yang diminati, mengisi formulir pemesanan, mengunggah bukti
        transfer, dan memantau status transaksi — semua dilakukan secara mandiri
        melalui \textit{browser} standar. Akses ke riwayat transaksi diamankan
        dengan verifikasi OTP berbasis nomor HP.

  \item \textbf{Antarmuka Administrasi (Kasir/Admin):} Dapat diakses hanya oleh
        staf internal showroom yang memiliki kredensial login. Jalur URL login
        sengaja disembunyikan dan hanya dapat ditemukan melalui
        \textit{keyboard shortcut} tersembunyi. Admin mengelola seluruh
        inventaris, memverifikasi transaksi, dan memantau metrik keuangan
        melalui \textit{dashboard} terintegrasi.
\end{enumerate}

\textit{Dashboard} admin menampilkan empat indikator KPI (\textit{Key Performance
Indicator}) yang diperbarui secara \textit{real-time}:

\begin{itemize}
  \item \textbf{Total Unit} --- jumlah keseluruhan unit kendaraan yang terdaftar
        dalam inventaris, tanpa memandang status ketersediaannya.
  \item \textbf{Unit Sold} --- jumlah unit yang statusnya telah berubah menjadi
        \texttt{Sold}, menandakan transaksi selesai dan unit telah dikonfirmasi
        serah terima.
  \item \textbf{Pending Verification} --- jumlah invoice yang memiliki bukti
        pembayaran terunggah dan menunggu persetujuan kasir.
  \item \textbf{Total Revenue} --- akumulasi nilai \texttt{total\_amount} dari
        seluruh invoice yang berstatus \texttt{Paid}, mencerminkan total
        pendapatan terkonfirmasi showroom.
\end{itemize}

% ----------------------------------------------------------------
\subsection*{2.2. Stakeholder dan User}
% ----------------------------------------------------------------

Pengguna yang berinteraksi dengan sistem ShowDrive terbagi menjadi dua kategori
dengan hak akses yang terpisah secara tegas:

\subsubsection*{1. Pelanggan (Customer)}

\textbf{Deskripsi:} Calon pembeli kendaraan dari kalangan umum yang mengakses
website ShowDrive melalui \textit{browser} tanpa memerlukan proses pendaftaran
akun. Identitas pelanggan dikenali sistem melalui nomor HP yang diverifikasi
dengan OTP setiap kali ingin mengakses riwayat transaksi.

\textbf{Hak Akses:}
\begin{itemize}
  \item Menelusuri katalog kendaraan dan melakukan pencarian berdasarkan merek,
        model, atau nomor VIN.
  \item Melihat halaman detail kendaraan beserta galeri foto dan spesifikasi
        teknis lengkap.
  \item Mengisi formulir pemesanan kendaraan dengan data diri (nama, nomor HP,
        NIK opsional, pilihan tipe pembayaran, dan tanggal inspeksi).
  \item Menerima dan menyimpan invoice digital dengan kode unik format
        \texttt{SD/YYYY/NNNN}.
  \item Meminta kode OTP dan melakukan verifikasi untuk mengakses
        \textit{dashboard} pelacakan transaksi pribadi.
  \item Mengunggah bukti transfer pembayaran (JPEG/PNG/WebP, maks.\ 2~MB).
  \item Membatalkan pemesanan secara mandiri selama status masih
        \texttt{Unpaid + Pending} disertai alasan pembatalan (min.\ 10 karakter).
  \item Mencetak kwitansi/invoice digital dari halaman pelacakan.
\end{itemize}

\subsubsection*{2. Administrator/Kasir (Admin)}

\textbf{Deskripsi:} Staf operasional internal showroom yang telah terdaftar
dalam sistem dan memiliki kredensial login. Kasir bertanggung jawab atas
verifikasi transaksi, pengelolaan inventaris, dan pemantauan kondisi bisnis
showroom secara keseluruhan.

\textbf{Hak Akses:}
\begin{itemize}
  \item Masuk ke panel administrasi melalui jalur URL tersembunyi yang diakses
        via \textit{keyboard shortcut} \texttt{Ctrl+Shift+A}.
  \item Memantau \textit{dashboard} metrik keuangan \textit{real-time}.
  \item Mengelola data kendaraan (CRUD: tambah, lihat, ubah, hapus) beserta
        galeri foto multi-gambar.
  \item Memverifikasi atau menolak bukti pembayaran yang diunggah pelanggan,
        disertai catatan penolakan yang dikirim sebagai notifikasi WhatsApp
        prefilled ke pelanggan.
  \item Melakukan \textit{amend} (perubahan tanggal inspeksi dan tipe pembayaran)
        pada reservasi yang belum dibayar.
  \item Mengonfirmasi serah terima fisik unit kendaraan setelah transaksi lunas,
        yang mengunci data transaksi secara permanen (\texttt{handed\_over\_at}).
  \item Membatalkan reservasi dari sisi admin dengan alasan wajib.
  \item Mengelola data gudang, akun kasir, profil perusahaan, dan pengaturan
        rekening bank/QRIS toko.
  \item Mengakses laporan keuangan ringkas.
\end{itemize}

% ----------------------------------------------------------------
\subsection*{2.3. Kebutuhan Fungsional}
% ----------------------------------------------------------------

Kebutuhan fungsional mendefinisikan layanan spesifik yang wajib disediakan oleh
sistem ShowDrive. Setiap fitur dirinci berdasarkan hak akses penggunanya.

\subsubsection*{A. Fitur pada Sisi Pelanggan}

\begin{enumerate}[label=\arabic*.]
  \item \textbf{Katalog Kendaraan dengan Pencarian:} Sistem menampilkan daftar
        unit kendaraan secara dinamis dengan informasi merek, model, tahun, harga
        (sebelum PPN), warna, dan status ketersediaan. Fitur pencarian
        \textit{full-text} mendukung kueri berdasarkan merek, model, dan nomor
        VIN (\textit{Vehicle Identification Number}).

  \item \textbf{Detail Kendaraan dan Galeri:} Halaman detail menampilkan
        spesifikasi teknis lengkap (mesin, transmisi, VIN, gudang penyimpanan)
        beserta galeri foto multi-gambar yang dapat dilihat satu per satu.

  \item \textbf{Formulir Pemesanan dengan Kalkulasi PPN:} Sistem menyediakan
        formulir pemesanan dengan validasi otomatis. Saat pelanggan memilih tipe
        pembayaran (DP atau lunas), sistem secara transparan menampilkan rincian:
        harga unit (subtotal), PPN 11\%, dan total yang harus dibayar. Untuk DP,
        sistem juga menampilkan nominal DP minimum berdasarkan persentase yang
        dikonfigurasi admin per unit.

  \item \textbf{Verifikasi Identitas Berbasis OTP:} Sebelum dapat mengakses
        riwayat transaksi, pelanggan harus memasukkan nomor HP dan memverifikasi
        kode OTP 4 digit yang di-\textit{generate} sistem. OTP berlaku selama
        5 menit, dibatasi 3 percobaan per sesi, dan di-\textit{hash} dengan
        Bcrypt sebelum disimpan di \textit{session} server.

  \item \textbf{Pengunggahan Bukti Pembayaran:} Sistem menerima unggahan gambar
        bukti transfer (JPEG, PNG, WebP; maks.\ 2~MB). Format SVG secara eksplisit
        diblokir untuk mencegah potensi serangan \textit{Cross-Site Scripting}
        (XSS) melalui file SVG berbahaya. Jika pelanggan mengunggah ulang bukti
        (misalnya karena ditolak admin), file lama dihapus secara aman melalui
        \textit{background job} setelah transaksi basis data berhasil.

  \item \textbf{Pelacakan Status dan Cetak Invoice:} Dari \textit{dashboard}
        pelacakan, pelanggan dapat memantau status pembayaran
        (\texttt{Unpaid} / \texttt{Pending Validation} / \texttt{Down Payment}
        / \texttt{Paid} / \texttt{Cancelled}) dan status pemesanan
        (\texttt{Pending} / \texttt{Approved} / \texttt{Rejected} /
        \texttt{Cancelled}), serta mencetak invoice digital.

  \item \textbf{Pembatalan Mandiri:} Pelanggan dapat membatalkan pemesanan
        sendiri selama status \texttt{payment\_status = Unpaid} dan
        \texttt{status = Pending}, dengan syarat memasukkan alasan pembatalan
        minimal 10 karakter. Status unit kendaraan otomatis dikembalikan ke
        \texttt{Available}.
\end{enumerate}

\subsubsection*{B. Fitur pada Sisi Administrator/Kasir}

\begin{enumerate}[label=\arabic*.]
  \item \textbf{Autentikasi Admin Berlapis:} Sistem menyediakan gerbang login
        dengan jalur URL tersembunyi (\texttt{/pintu-akses-masuk-showdrive}) yang
        hanya dapat ditemukan melalui \textit{shortcut} keyboard. \textit{Rate
        limiter} membatasi percobaan login maksimal 5 kali per menit per alamat
        IP untuk mencegah serangan \textit{brute force}.

  \item \textbf{Dashboard Metrik Real-Time:} \textit{Dashboard} menampilkan 4
        KPI utama yang diperbarui dinamis, serta tabel ringkasan transaksi
        terbaru.

  \item \textbf{Manajemen Kendaraan (CRUD):} Admin dapat menambah unit kendaraan
        baru dengan data lengkap (merek, model, VIN, tahun, harga, persentase DP,
        mesin, transmisi, warna, gudang, foto utama), mengedit, menghapus (dengan
        proteksi \texttt{ON DELETE RESTRICT} jika unit masih terikat invoice),
        serta mengelola galeri foto tambahan per unit.

  \item \textbf{Verifikasi Pembayaran via AJAX:} Kasir dapat menyetujui atau
        menolak bukti pembayaran secara asinkron. Persetujuan secara otomatis
        mengubah \texttt{payment\_status} invoice, mencatat \texttt{cashier\_id}
        yang menyetujui, dan memperbarui status unit kendaraan. Penolakan
        mengharuskan kasir mengisi catatan alasan minimal 10 karakter. Setelah
        aksi berhasil, sistem menyiapkan \textit{link} WhatsApp prefilled dengan
        pesan notifikasi yang dipersonalisasi untuk dikirim kasir ke pelanggan
        secara manual.

  \item \textbf{Amend Reservasi:} Kasir dapat mengubah tanggal inspeksi
        (terbatas dalam rentang +1 hingga +7 hari dari hari ini) dan tipe
        pembayaran (DP atau lunas) pada reservasi yang masih berstatus belum
        bayar.

  \item \textbf{Konfirmasi Serah Terima Unit (\textit{Handover}):} Setelah
        pembayaran lunas terverifikasi, kasir mengonfirmasi serah terima fisik
        unit kendaraan. Sistem mencatat \texttt{handed\_over\_at} (waktu serah
        terima) dan mengunci data transaksi dari modifikasi lebih lanjut,
        menjamin keutuhan rekaman historis.

  \item \textbf{Edit Data Pelanggan:} Kasir dapat mengoreksi nama dan NIK
        pelanggan pada invoice yang belum di-\textit{handover} dan belum
        dibatalkan, untuk menangani kesalahan penginputan oleh pelanggan.

  \item \textbf{Cetak Invoice Dinamis:} Sistem menyediakan dua jenis invoice
        yang dapat dicetak: invoice DP (menampilkan nominal uang muka) dan
        kwitansi lunas (menampilkan total pembayaran lengkap dengan PPN).
\end{enumerate}

% ----------------------------------------------------------------
\subsection*{2.4. Kebutuhan Non-Fungsional}
% ----------------------------------------------------------------

Kebutuhan non-fungsional mendefinisikan batasan kualitas, performa, dan keamanan
yang harus dipenuhi sistem ShowDrive agar dapat dioperasikan secara optimal dalam
lingkungan produksi.

\begin{enumerate}[label=\alph*.]
  \item \textbf{Kinerja (\textit{Performance}):} Halaman katalog menggunakan
        proyeksi kolom selektif (\texttt{select()} pada \texttt{Query Builder})
        untuk menghindari pengambilan data yang tidak diperlukan ke memori. Proses
        verifikasi pembayaran via AJAX didesain merespons dalam waktu kurang dari
        1 detik pada kondisi jaringan lokal yang stabil.

  \item \textbf{Kemudahan Penggunaan (\textit{Usability}):} Antarmuka publik
        dirancang agar dapat digunakan tanpa panduan, dengan alur transaksi
        linear yang jelas. Antarmuka admin menggunakan dropdown aksi per baris
        (\textit{action dropdown}) yang menutup otomatis saat area luar di-klik,
        mengurangi kepadatan visual pada halaman daftar transaksi.

  \item \textbf{Keamanan (\textit{Security}):} Sistem mengimplementasikan empat
        lapisan keamanan yang bekerja secara berlapis:
        \begin{itemize}
          \item \textit{Security by obscurity}: URL login admin disembunyikan dari
                publik.
          \item \textit{Rate limiting}: pembatasan permintaan OTP (5x/10 menit),
                percobaan login (5x/menit), dan pembuatan booking (10x/menit) per
                alamat IP.
          \item \textit{Session-based OTP}: OTP di-\textit{hash} Bcrypt sebelum
                disimpan dan diverifikasi menggunakan \texttt{Hash::check()},
                mencegah pembandingan nilai \textit{plaintext} meski sesi
                diinspeksi.
          \item \textit{Database-level integrity}: \textit{constraint}
                \texttt{ON DELETE} pada tingkat basis data sebagai jaring pengaman
                terakhir, independen dari lapisan validasi aplikasi.
        \end{itemize}

  \item \textbf{Keandalan (\textit{Reliability}):} Seluruh operasi multi-tabel
        (pembuatan booking, upload bukti, verifikasi pembayaran, pembatalan)
        dibungkus dalam \texttt{DB::transaction()}, memastikan tidak ada kondisi
        data yang inkonsisten akibat kegagalan di tengah operasi. File bukti
        pembayaran lama hanya dihapus melalui \textit{background job} setelah
        transaksi basis data berhasil di-\textit{commit}.

  \item \textbf{Skalabilitas:} Normalisasi basis data ke dalam 6 tabel dengan
        \textit{Foreign Key} terstruktur memastikan sistem dapat dikembangkan
        dengan fitur baru (multi-cabang, \textit{payment gateway}, notifikasi
        otomatis) tanpa perombakan skema yang mendasar.
\end{enumerate}

% ----------------------------------------------------------------
\subsection*{2.5. Keunggulan Sistem}
% ----------------------------------------------------------------

Bagian ini menjabarkan keunggulan teknis dan bisnis ShowDrive secara mendalam,
dirancang sebagai referensi bagi pembaca yang ingin memahami fondasi arsitektur
sistem sebelum masuk ke tahap perancangan dan implementasi.

\subsubsection*{a. Mekanisme Transaksi Atomik (Prinsip ACID)}

Tantangan terbesar dalam sistem transaksi multi-tabel adalah menjamin bahwa
semua operasi yang saling bergantung berhasil secara keseluruhan, atau tidak ada
yang berhasil sama sekali. Pada ShowDrive, proses pemesanan kendaraan secara
bersamaan memengaruhi tiga tabel: \texttt{customers} (upsert data pelanggan),
\texttt{invoices} (buat rekaman transaksi baru dengan kalkulasi PPN), dan
\texttt{items} (ubah status dari \texttt{Available} menjadi \texttt{Invoiced}).

Sistem mengimplementasikan \texttt{DB::transaction()} untuk membungkus ketiga
operasi tersebut. Jika salah satu operasi gagal (misalnya karena NIK yang
dimasukkan sudah terdaftar di akun lain dan memicu
\texttt{UniqueConstraintViolationException}), seluruh transaksi di-\textit{rollback}
otomatis dan tidak ada perubahan parsial yang tersimpan di basis data. Pendekatan
ini mengimplementasikan prinsip \textit{Atomicity} dari ACID
(Silberschatz et al., 2020) sebagai fondasi keandalan sistem.

\begin{infobox}
\textbf{Implikasi Bisnis:} Tidak akan pernah terjadi kondisi di mana invoice
sudah terbuat namun status unit masih \texttt{Available} (sehingga bisa dipesan
dua kali), atau status unit berubah menjadi \texttt{Invoiced} namun tidak ada
invoice yang terbuat di sistem.
\end{infobox}

\subsubsection*{b. Integritas Referensial Berlapis dengan Tiga Kebijakan
\textit{ON DELETE}}

Berbeda dari sistem yang hanya mengandalkan validasi di lapisan aplikasi,
ShowDrive menerapkan integritas referensial langsung di tingkat basis data MySQL
melalui tiga kebijakan \textit{constraint ON DELETE} yang berbeda, disesuaikan
dengan logika bisnis masing-masing relasi:

\begin{itemize}
  \item \textbf{\texttt{ON DELETE RESTRICT}} pada relasi \texttt{items}
        $\rightarrow$ \texttt{invoices}: Basis data akan \textit{menolak} operasi
        \texttt{DELETE} pada tabel \texttt{items} jika unit kendaraan tersebut
        masih memiliki satu atau lebih rekaman invoice. Ini adalah proteksi
        keandalan paling kritis — menjamin bahwa tidak ada aset kendaraan bernilai
        miliaran rupiah yang terhapus secara tidak sengaja oleh kasir selama masih
        ada transaksi aktif yang mengacu padanya.

  \item \textbf{\texttt{ON DELETE CASCADE}} pada relasi \texttt{customers}
        $\rightarrow$ \texttt{invoices}: Jika data pelanggan dihapus dari sistem,
        seluruh riwayat invoice milik pelanggan tersebut juga ikut dihapus secara
        atomik. Kebijakan ini mendukung kebutuhan pengelolaan data pelanggan yang
        sudah tidak aktif tanpa meninggalkan data \textit{orphan} (rekaman tanpa
        induk) di tabel invoices.

  \item \textbf{\texttt{ON DELETE SET NULL}} pada relasi \texttt{cashiers}
        $\rightarrow$ \texttt{invoices}: Jika akun kasir dihapus dari sistem,
        kolom \texttt{cashier\_id} pada invoice yang pernah diverifikasi oleh
        kasir tersebut diset menjadi \texttt{NULL}. Data historis transaksi tetap
        utuh dan terjaga, tidak ikut terhapus meski kasir yang bersangkutan
        sudah tidak aktif.
\end{itemize}

\subsubsection*{c. Sistem OTP Berlapis untuk Perlindungan Data Transaksi}

Sistem pelacakan transaksi menggunakan OTP sebagai pengganti sistem login
konvensional, memberikan keseimbangan antara keamanan dan kemudahan penggunaan.
Alur kerja OTP yang diimplementasikan:

\begin{enumerate}[label=\arabic*.]
  \item Pelanggan memasukkan nomor HP; sistem memvalidasi format dan memastikan
        nomor HP terdaftar beserta riwayat invoicenya.
  \item Sistem men-\textit{generate} OTP acak 4 digit
        (\texttt{random\_int(1000, 9999)}), meng-\textit{hash}-nya dengan Bcrypt
        (\texttt{bcrypt((string) \$otp)}), dan menyimpannya di \textit{session}
        server beserta \textit{timestamp} kedaluwarsa 5 menit.
  \item OTP kemudian ditampilkan sebagai simulasi (mode pengembangan) atau
        dikirim via saluran eksternal (mode produksi).
  \item Verifikasi menggunakan \texttt{Hash::check()} — bukan perbandingan
        \textit{string} langsung — sehingga nilai OTP yang sebenarnya tidak
        pernah dapat diekstrak dari \textit{session} meski sesi diinspeksi.
  \item Setelah 3 kali percobaan gagal, OTP hangus dan sesi dihapus.
        \textit{Rate limiter} pada level \textit{routing} menambah lapisan
        perlindungan terhadap serangan multi-sesi.
\end{enumerate}

\subsubsection*{d. Kode Invoice Sequential yang Aman dari \textit{Race Condition}}

Format kode invoice ShowDrive (\texttt{SD/2026/0001}, \texttt{SD/2026/0002},
dst.) menggunakan algoritma \textit{sequential numbering} yang dirancang aman
dari \textit{race condition} — kondisi di mana dua pengguna melakukan booking
bersamaan dan berpotensi mendapatkan nomor yang sama.

Keamanan ini dicapai melalui kombinasi \texttt{DB::transaction()} di lapisan
luar dan \texttt{lockForUpdate()} saat mengambil nomor invoice terakhir. Dengan
\textit{row-level lock} ini, request pertama yang mendapatkan \textit{lock} akan
menyelesaikan operasinya dan meng-\textit{commit} nomor baru ke basis data
sebelum request kedua dapat membaca nilai nomor terakhir, sehingga duplikasi
nomor invoice secara teknis tidak mungkin terjadi.

\newpage

% ================================================================
%  BAB 3  PERANCANGAN SISTEM
% ================================================================
\section*{BAB 3 PERANCANGAN SISTEM}
\addcontentsline{toc}{section}{BAB 3 PERANCANGAN SISTEM}

% ----------------------------------------------------------------
\subsection*{3.1. Arsitektur Sistem}
% ----------------------------------------------------------------

Aplikasi ShowDrive dibangun menggunakan arsitektur tiga lapis
(\textit{three-tier architecture}) yang memisahkan tanggung jawab antara
tampilan, logika bisnis, dan penyimpanan data secara modular. Pemisahan ini
memastikan bahwa perubahan pada satu lapisan tidak memengaruhi lapisan lainnya,
sehingga sistem lebih mudah dikembangkan, diuji, dan dipelihara.

\begin{figure}[H]
  \centering
  \begin{tikzpicture}[node distance=0.45cm and 0.5cm]

    %% ---- LAYER BOXES ----
    \node[bigbox, fill=blue!10, minimum width=7.5cm, minimum height=1.1cm]
      (fe) {PRESENTATION LAYER\\[2pt]{\footnotesize\normalfont Frontend}};

    \node[bigbox, fill=green!10, minimum width=7.5cm, minimum height=1.1cm,
          below=0.7cm of fe]
      (be) {APPLICATION LAYER\\[2pt]{\footnotesize\normalfont Backend}};

    \node[bigbox, fill=orange!10, minimum width=7.5cm, minimum height=1.1cm,
          below=0.7cm of be]
      (db) {DATABASE LAYER\\[2pt]{\footnotesize\normalfont Penyimpanan Data}};

    %% ---- TEKNOLOGI dalam setiap layer ----
    % Frontend
    \node[box, fill=blue!20, below=0.1cm of fe.north,
          yshift=-0.45cm, xshift=-1.6cm, minimum width=1.8cm]
      (blade) {Blade\\Template};
    \node[box, fill=blue!20, right=0.25cm of blade, minimum width=1.8cm]
      (tailwind) {Tailwind\\CSS};
    \node[box, fill=blue!20, right=0.25cm of tailwind, minimum width=1.8cm]
      (js) {JavaScript\\AJAX};

    % Backend
    \node[box, fill=green!20, below=0.1cm of be.north,
          yshift=-0.45cm, xshift=-1.6cm, minimum width=1.8cm]
      (laravel) {Laravel 12\\PHP 8.2};
    \node[box, fill=green!20, right=0.25cm of laravel, minimum width=1.8cm]
      (ctrl) {Controller\\+ Middleware};
    \node[box, fill=green!20, right=0.25cm of ctrl, minimum width=1.8cm]
      (orm) {Eloquent\\ORM};

    % Database
    \node[box, fill=orange!20, below=0.1cm of db.north,
          yshift=-0.45cm, xshift=-1.6cm, minimum width=1.8cm]
      (mysql) {MySQL 8.0\\RDBMS};
    \node[box, fill=orange!20, right=0.25cm of mysql, minimum width=1.8cm]
      (schema) {6 Tabel\\Relasional};
    \node[box, fill=orange!20, right=0.25cm of schema, minimum width=1.8cm]
      (fk) {Foreign Key\\Constraints};

    %% ---- PANAH ANTAR LAYER ----
    \draw[barr] ([yshift=-0.1cm]fe.south) -- ([yshift=0.1cm]be.north)
      node[midway, right=0.2cm, font=\scriptsize] {HTTP Request/Response};
    \draw[barr] ([yshift=-0.1cm]be.south) -- ([yshift=0.1cm]db.north)
      node[midway, right=0.2cm, font=\scriptsize] {SQL / Eloquent ORM};

    %% ---- AKTOR DI KIRI ----
    \node[box, fill=gray!15, left=0.9cm of fe, minimum width=1.4cm,
          minimum height=0.65cm, text width=1.3cm, font=\scriptsize]
      (customer) {Pelanggan\\(Browser)};
    \node[box, fill=gray!15, left=0.9cm of be, minimum width=1.4cm,
          minimum height=0.65cm, text width=1.3cm, font=\scriptsize]
      (admin) {Admin /\\Kasir};

    \draw[arr] (customer.east) -- (fe.west);
    \draw[arr] (admin.east)    -- (be.west);

  \end{tikzpicture}
  \caption{Arsitektur Sistem ShowDrive (\textit{Three-Tier Architecture})}
  \label{fig:arsitektur}
\end{figure}

Ketiga lapisan arsitektur ShowDrive bekerja sebagai berikut:

\begin{enumerate}[label=\arabic*.]
  \item \textbf{Lapisan Presentasi / \textit{Frontend} (View):} Merupakan lapisan
        antarmuka yang berinteraksi langsung dengan pengguna. Dibangun menggunakan
        Blade Template Engine Laravel (\texttt{.blade.php}) sebagai mesin
        \textit{rendering} HTML dinamis, Tailwind CSS untuk desain visual responsif
        dengan tema \textit{Premium Dark Mode}, serta JavaScript murni (tanpa
        \textit{framework} tambahan) untuk operasi asinkron (AJAX/Fetch API) seperti
        verifikasi pembayaran, amend reservasi, dan konfirmasi serah terima unit
        tanpa \textit{page reload}.

  \item \textbf{Lapisan Aplikasi / \textit{Backend} (Controller \& Logic):}
        Dibangun menggunakan Laravel 12 dengan PHP 8.2. Lapisan ini menangani
        seluruh logika bisnis melalui \textit{Controller}, memvalidasi input
        pengguna melalui \textit{Form Request} (mis.\ \texttt{StoreBookingRequest}),
        memproses kalkulasi finansial (PPN, DP), mengelola \textit{session} OTP,
        dan mengeksekusi operasi basis data melalui Eloquent ORM. \textit{Middleware}
        digunakan untuk membatasi akses halaman admin hanya pada pengguna
        terautentikasi, dan \textit{Rate Limiter} melindungi \textit{endpoint}
        kritis dari penyalahgunaan.

  \item \textbf{Lapisan Basis Data (Model \& Database):} Menggunakan MySQL 8.0
        sebagai RDBMS dengan skema enam tabel relasional yang telah dinormalisasi.
        Eloquent ORM Laravel mengelola interaksi antara objek PHP dan tabel basis
        data, sementara mekanisme \textit{migration} menjamin konsistensi skema
        di semua lingkungan pengembangan.
\end{enumerate}

% ----------------------------------------------------------------
\subsection*{3.2. Workflow Sistem}
% ----------------------------------------------------------------

Alur kerja ShowDrive terbagi menjadi dua jalur utama berdasarkan peran pengguna.
Kedua jalur ini saling berinteraksi melalui satu basis data yang sama.

\subsubsection*{A. Workflow Pelanggan (\textit{Customer Workflow})}

\begin{figure}[H]
  \centering
  \begin{tikzpicture}[node distance=0.38cm and 1.2cm]

    % --- KOLOM KIRI: node ---
    \node[term] (start) {MULAI};

    \node[proc, below=of start]
      (katalog) {1. Buka Halaman Katalog\\(\texttt{GET /})};

    \node[proc, below=of katalog]
      (detail) {2. Lihat Detail Kendaraan\\(\texttt{GET /car/\{id\}})};

    \node[proc, below=of detail]
      (form) {3. Isi Formulir Pemesanan\\(nama, HP, NIK, tipe bayar)};

    \node[proc, below=of form]
      (atomik) {4. Sistem Proses Atomik\\(\texttt{DB::transaction()})};

    \node[proc, below=of atomik]
      (otp) {5. Minta \& Verifikasi OTP\\(\texttt{POST /track/otp})};

    \node[dec, below=of otp]
      (otpok) {OTP Valid?};

    \node[proc, below=of otpok]
      (upload) {6. Unggah Bukti Transfer\\(\texttt{POST /.../upload-proof})};

    \node[proc, below=of upload]
      (lacak) {7. Pantau Status Invoice\\(\texttt{GET /track})};

    \node[proc, below=of lacak]
      (cetak) {8. Cetak Invoice / Kwitansi\\(\texttt{GET /invoice/\{id\}})};

    \node[term, below=of cetak] (end) {SELESAI};

    % --- KETERANGAN SAMPINGAN ---
    \node[box, fill=yellow!20, right=1.4cm of atomik,
          text width=2.5cm, minimum height=1.1cm, font=\scriptsize]
      (atombox) {Upsert customers\\Buat invoices\\Update items $\to$ Invoiced};

    \node[box, fill=red!15, right=1.4cm of otpok,
          text width=2.2cm, minimum height=0.65cm, font=\scriptsize]
      (otpfail) {Sisa percobaan\\berkurang (maks 3)};

    \node[box, fill=gray!15, right=1.4cm of upload,
          text width=2.2cm, minimum height=0.8cm, font=\scriptsize]
      (statusbox) {Status $\to$\\Pending Validation};

    % --- PANAH UTAMA ---
    \foreach \a/\b in {start/katalog, katalog/detail, detail/form,
                       form/atomik, atomik/otp, otp/otpok,
                       upload/lacak, lacak/cetak, cetak/end}
      \draw[arr] (\a) -- (\b);

    \draw[arr] (otpok.south) -- node[right,font=\scriptsize]{Ya} (upload.north);
    \draw[arr] (otpok.east)  -| node[above,font=\scriptsize]{Tidak}
               ([xshift=0.3cm]otpfail.north) -- (otpfail.north);
    \draw[darr] (otpfail.west) -- ++(-0.5,0) |- (otp.east);
    \draw[darr] (atomik.east) -- (atombox.west);
    \draw[darr] (upload.east) -- (statusbox.west);

  \end{tikzpicture}
  \caption{Workflow Pelanggan ShowDrive}
  \label{fig:workflow-pelanggan}
\end{figure}

Berikut adalah narasi lengkap alur pelanggan dari awal hingga selesai:

\begin{enumerate}[label=\arabic*.]
  \item \textbf{Penelusuran Katalog:} Pelanggan membuka halaman beranda
        (\texttt{GET /}) yang menampilkan daftar unit kendaraan berstatus
        \texttt{Available} dalam tampilan grid. Pelanggan dapat menggunakan fitur
        pencarian \textit{real-time} untuk memfilter unit berdasarkan merek, model,
        atau nomor VIN.

  \item \textbf{Detail Kendaraan:} Pelanggan mengklik unit yang diminati dan
        diarahkan ke halaman detail (\texttt{GET /car/\{id\}}) yang menampilkan
        spesifikasi teknis lengkap, lokasi gudang penyimpanan, serta galeri foto
        tambahan.

  \item \textbf{Pengisian Formulir Pemesanan:} Pelanggan mengisi formulir
        pemesanan dengan data: nama lengkap, nomor HP aktif, NIK (opsional),
        tanggal inspeksi yang diinginkan, dan tipe pembayaran (DP atau lunas).
        Sistem secara otomatis menampilkan rincian kalkulasi harga termasuk PPN
        11\% sebelum formulir dikirim.

  \item \textbf{Pemrosesan Booking Atomik:} Saat pelanggan menekan tombol
        \textit{submit}, sistem (\texttt{POST /booking}) memproses tiga operasi
        secara atomik dalam satu transaksi basis data:
        \begin{itemize}
          \item \textbf{Upsert pelanggan:} Mencari data pelanggan berdasarkan
                nomor HP. Jika belum ada, data baru dibuat. Jika sudah ada dan
                NIK belum terisi, NIK diperbarui.
          \item \textbf{Buat invoice:} Rekaman transaksi baru dibuat dengan
                kode unik \texttt{SD/YYYY/NNNN}, kalkulasi subtotal, PPN 11\%,
                dan total akhir; status awal \texttt{payment\_status = Unpaid},
                \texttt{status = Pending}.
          \item \textbf{Kunci unit:} Status unit di tabel \texttt{items} diubah
                dari \texttt{Available} menjadi \texttt{Invoiced} (tampil sebagai
                \texttt{Booked} di antarmuka publik).
        \end{itemize}

  \item \textbf{Verifikasi OTP untuk Akses Riwayat:} Untuk mengakses
        \textit{dashboard} pelacakan, pelanggan memasukkan nomor HP lalu meminta
        OTP (\texttt{POST /track/otp}). Setelah menerima OTP, pelanggan
        memasukkanya untuk diverifikasi (\texttt{POST /track/verify}). Sesi
        terverifikasi disimpan di \textit{session} server.

  \item \textbf{Unggah Bukti Pembayaran:} Dari \textit{dashboard} pelacakan
        (\texttt{GET /track}), pelanggan mengunggah foto bukti transfer
        (\texttt{POST /booking/\{id\}/upload-proof}). Status invoice otomatis
        berubah menjadi \texttt{Pending Validation}. Jika pelanggan mengunggah
        ulang (mengganti bukti), file lama dihapus aman via
        \texttt{DeleteOldImageJob}.

  \item \textbf{Pemantauan Status:} Pelanggan memantau perubahan status pembayaran
        dan status pemesanan dari halaman yang sama. Status akan berubah sesuai
        aksi kasir (disetujui, ditolak, atau dibatalkan).

  \item \textbf{Cetak Invoice:} Setelah pembayaran terverifikasi, pelanggan
        dapat mencetak kwitansi digital melalui \texttt{GET /invoice/\{id\}} yang
        menampilkan detail lengkap transaksi dalam format siap cetak.

  \item \textbf{Pembatalan Mandiri (Opsional):} Selama status masih
        \texttt{Unpaid + Pending}, pelanggan dapat membatalkan pemesanan sendiri
        (\texttt{POST /booking/\{id\}/cancel}) dengan alasan yang wajib diisi.
        Status unit otomatis dikembalikan ke \texttt{Available}.
\end{enumerate}

\subsubsection*{B. Workflow Administrator/Kasir (\textit{Admin Workflow})}

\begin{figure}[H]
  \centering
  \begin{tikzpicture}[node distance=0.38cm and 1.2cm]

    % --- KOLOM KIRI: node ---
    \node[term] (start) {MULAI};

    \node[proc, below=of start]
      (login) {1. Login via Shortcut\\(\texttt{Ctrl+Shift+A})};

    \node[proc, below=of login]
      (dash) {2. Pantau Dashboard\\4 KPI Real-Time};

    \node[proc, below=of dash]
      (list) {3. Buka Halaman\\Validasi Transaksi};

    \node[dec, below=of list]
      (hasbukti) {Ada Bukti\\Transfer?};

    \node[proc, below=of hasbukti]
      (tinjau) {4. Tinjau Gambar\\Bukti Pembayaran};

    \node[dec, below=of tinjau]
      (valid) {Bukti\\Valid?};

    \node[proc, below=of valid]
      (approve) {5. Klik Sahkan Lunas\\(AJAX approve)};

    \node[proc, below=of approve]
      (notif) {6. Kirim Notif WA\\(prefilled link)};

    \node[proc, below=of notif]
      (handover) {7. Konfirmasi\\Serah Terima Unit};

    \node[term, below=of handover] (end) {DATA TERKUNCI};

    % --- KETERANGAN SAMPINGAN ---
    \node[box, fill=red!15, right=1.3cm of valid,
          text width=2.2cm, minimum height=0.8cm, font=\scriptsize]
      (tolak) {Isi catatan\\penolakan (min 10\\karakter) $\to$ Tolak};

    \node[box, fill=green!15, right=1.3cm of approve,
          text width=2.4cm, minimum height=1.0cm, font=\scriptsize]
      (effect) {invoice $\to$ Paid\\cashier\_id dicatat\\items $\to$ Sold};

    \node[box, fill=gray!15, right=1.3cm of handover,
          text width=2.4cm, minimum height=0.8cm, font=\scriptsize]
      (lock) {handed\_over\_at\\dicatat. Semua\\aksi terkunci.};

    \node[box, fill=yellow!20, right=1.3cm of list,
          text width=2.4cm, minimum height=0.8cm, font=\scriptsize]
      (amend) {Opsional: Amend\\tanggal inspeksi\\/ tipe bayar};

    % --- PANAH UTAMA ---
    \foreach \a/\b in {start/login, login/dash, dash/list,
                       list/hasbukti, tinjau/valid,
                       approve/notif, notif/handover, handover/end}
      \draw[arr] (\a) -- (\b);

    \draw[arr] (hasbukti.south) -- node[right,font=\scriptsize]{Ya} (tinjau.north);
    \draw[darr] (hasbukti.west) -- ++(-0.6,0)
      node[left,font=\scriptsize]{Tidak\\(tunggu)};

    \draw[arr] (valid.south) -- node[right,font=\scriptsize]{Ya} (approve.north);
    \draw[arr] (valid.east)  -| node[above,font=\scriptsize]{Tidak}
               ([xshift=0.3cm]tolak.north) -- (tolak.north);
    \draw[darr] (tolak.south) -- ++(0,-0.3) -|
               ([xshift=1.3cm]list.east) -- (list.east);

    \draw[darr] (approve.east) -- (effect.west);
    \draw[darr] (handover.east) -- (lock.west);
    \draw[darr] (list.east) -- (amend.west);

  \end{tikzpicture}
  \caption{Workflow Admin/Kasir ShowDrive}
  \label{fig:workflow-admin}
\end{figure}

\begin{enumerate}[label=\arabic*.]
  \item \textbf{Autentikasi Admin:} Kasir menekan \texttt{Ctrl+Shift+A} di
        halaman manapun untuk diarahkan ke \texttt{/pintu-masuk}, yang kemudian
        me-\textit{redirect} ke halaman login tersembunyi
        \texttt{/pintu-akses-masuk-showdrive}. Kasir memasukkan \textit{username}
        dan \textit{password}; sistem memverifikasi dengan \textit{rate limiter}
        5 percobaan per menit per IP.

  \item \textbf{Pemantauan Dashboard:} Setelah login, kasir diarahkan ke
        \texttt{/admin/dashboard} yang menampilkan 4 KPI utama dan tabel
        ringkasan transaksi terbaru secara \textit{real-time}.

  \item \textbf{Pengelolaan Transaksi:} Kasir membuka halaman
        \texttt{/admin/invoices} yang menampilkan semua invoice dengan filter
        status (\textit{Pending}, \textit{Down Payment}, \textit{Lunas},
        \textit{Dibatalkan}). Setiap baris memiliki \textit{action dropdown}
        dengan opsi aksi kontekstual.

  \item \textbf{Verifikasi atau Penolakan Pembayaran:} Kasir meninjau gambar
        bukti transfer yang diunggah pelanggan (klik \textit{``Lihat''} untuk
        membuka gambar di tab baru). Jika valid, kasir memilih \textit{``Setujui
        \& Sahkan''} (AJAX \texttt{POST /admin/booking/\{id\}/approve-ajax}).
        Jika tidak valid, kasir memilih \textit{``Tolak Bukti''} dan mengisi
        catatan alasan penolakan.

  \item \textbf{Notifikasi WhatsApp Prefilled:} Setelah verifikasi atau
        penolakan berhasil, sistem secara otomatis membuka \textit{link}
        WhatsApp prefilled dengan pesan notifikasi yang dipersonalisasi (nama
        pelanggan, nama unit, kode invoice, dan — khusus penolakan — alasan
        penolakan) untuk dikirim kasir ke pelanggan.

  \item \textbf{Amend Reservasi (Opsional):} Untuk reservasi yang belum dibayar,
        kasir dapat mengubah tanggal inspeksi atau tipe pembayaran melalui modal
        \textit{amend} (AJAX \texttt{POST /admin/booking/\{id\}/amend-ajax}).

  \item \textbf{Konfirmasi Serah Terima (\textit{Handover}):} Setelah invoice
        berstatus \texttt{Paid}, kasir mengonfirmasi serah terima fisik unit
        (AJAX \texttt{POST /admin/booking/\{id\}/handover-ajax}). Sistem mencatat
        waktu serah terima di kolom \texttt{handed\_over\_at} dan mengunci seluruh
        data transaksi dari modifikasi lebih lanjut (\textit{``Data Terkunci''}).
\end{enumerate}

% ----------------------------------------------------------------
\subsection*{3.3. Siklus Hidup Status Transaksi}
% ----------------------------------------------------------------

Pemahaman atas siklus hidup status transaksi sangat penting bagi pembaca yang
ingin memahami logika bisnis ShowDrive secara menyeluruh. Setiap invoice memiliki
\textbf{dua dimensi status} yang bergerak secara independen namun saling
berkaitan: \texttt{payment\_status} dan \texttt{status} (inspeksi/pemesanan).

Tabel \ref{tab:lifecycle} merangkum seluruh kemungkinan transisi status beserta
kondisi pemicunya:

\begin{table}[H]
  \centering
  \small
  \renewcommand{\arraystretch}{1.4}
  \caption{Siklus Hidup Status Invoice ShowDrive}
  \label{tab:lifecycle}
  \begin{tabularx}{\textwidth}{|l|l|X|}
    \hline
    \textbf{payment\_status} & \textbf{status} & \textbf{Kondisi / Pemicu} \\
    \hline
    \texttt{Unpaid}          & \texttt{Pending}   & Status awal saat booking dibuat. \\
    \hline
    \texttt{Pending Validation} & \texttt{Pending} & Pelanggan mengunggah bukti transfer. \\
    \hline
    \texttt{Down Payment}    & \texttt{Approved}  & Kasir menyetujui bukti DP. Unit tetap
                                                     \texttt{Invoiced}. \\
    \hline
    \texttt{Pending Validation} & \texttt{Approved} & Pelanggan mengunggah bukti pelunasan
                                                      (setelah DP disetujui). \\
    \hline
    \texttt{Paid}            & \texttt{Approved}  & Kasir menyetujui bukti pelunasan. Unit
                                                     berubah menjadi \texttt{Sold}. \\
    \hline
    \texttt{Unpaid}          & \texttt{Rejected}  & Kasir menolak bukti bayar. Unit kembali
                                                     ke \texttt{Available}. \\
    \hline
    \texttt{Cancelled}       & \texttt{Cancelled} & Dibatalkan oleh pelanggan atau admin.
                                                     Unit kembali ke \texttt{Available}. \\
    \hline
  \end{tabularx}
\end{table}

Status unit kendaraan di tabel \texttt{items} juga memiliki siklus sendiri yang
bergerak seiring status invoice:

\begin{center}
\texttt{Available} $\xrightarrow{\text{booking dibuat}}$
\texttt{Invoiced} $\xrightarrow{\text{pelunasan disetujui}}$
\texttt{Sold}
\end{center}

\vspace{0.3cm}
Jika invoice dibatalkan atau bukti ditolak dari kondisi \texttt{Invoiced}, unit
kembali ke \texttt{Available}. Status \texttt{Sold} bersifat permanen setelah
\textit{handover} dikonfirmasi.

% ----------------------------------------------------------------
\subsection*{3.4. Class Diagram}
% ----------------------------------------------------------------

Class Diagram berikut menggambarkan struktur kelas model Eloquent yang
digunakan dalam \textit{framework} Laravel untuk mengatur logika bisnis dan
relasi antar entitas ShowDrive:

\begin{figure}[H]
  \centering
  \begin{tikzpicture}[node distance=0.5cm and 0.6cm]

    % ---- Definisi style kelas UML lokal ----
    \tikzstyle{clshead}=[rectangle, draw=black!80, fill=blue!15,
      minimum width=2.8cm, minimum height=0.55cm,
      font=\scriptsize\bfseries, text centered]
    \tikzstyle{clsbody}=[rectangle, draw=black!80, fill=blue!4,
      minimum width=2.8cm, font=\scriptsize,
      text width=2.7cm, align=left, inner sep=3pt]

    % ============ CLASS COMPANY ============
    \node[clshead] (ch) at (0,0) {Company};
    \node[clsbody, below=0cm of ch] (cb)
      {id, name, tax\_id\\phone, address\\bank\_name, qris\_image};

    % ============ CLASS WAREHOUSE ============
    \node[clshead, right=1.2cm of ch] (wh) {Warehouse};
    \node[clsbody, below=0cm of wh] (wb)
      {id, company\_id (FK)\\name, location};

    % ============ CLASS CASHIER ============
    \node[clshead, right=1.2cm of wh] (kh) {Cashier};
    \node[clsbody, below=0cm of kh] (kb)
      {id, company\_id (FK)\\name, username\\password (bcrypt)};

    % ============ CLASS ITEM ============
    \node[clshead, below=2.4cm of wh] (ih) {Item};
    \node[clsbody, below=0cm of ih] (ib)
      {id, warehouse\_id (FK)\\brand, model, vin\\year, price, dp\_pct\\status, image\_url};

    % ============ CLASS CUSTOMER ============
    \node[clshead, left=1.2cm of ih] (cuh) {Customer};
    \node[clsbody, below=0cm of cuh] (cub)
      {id, name\\phone (UNIQUE)\\nik (UNIQUE, NULL)};

    % ============ CLASS INVOICE ============
    \node[clshead, right=1.2cm of ih] (invh) {Invoice};
    \node[clsbody, below=0cm of invh] (invb)
      {id, invoice\_code\\customer\_id (FK)\\item\_id (FK)\\cashier\_id (FK, NULL)\\total\_amount, tax\\payment\_status\\status, handed\_over\_at};

    % ============ CLASS ITEMIMAGE ============
    \node[clshead, below=2.2cm of ih] (imgh) {ItemImage};
    \node[clsbody, below=0cm of imgh] (imgb)
      {id, item\_id (FK)\\image\_path};

    % ---- RELASI ----
    % Company 1:N Warehouse
    \draw[arr] (cb.south) -- ++(0,-0.3)
      -| node[above, font=\scriptsize, pos=0.25]{1:N} (wh.south);

    % Company 1:N Cashier
    \draw[arr] (cb.south) -- ++(0,-0.3)
      -| node[above, font=\scriptsize, pos=0.85]{1:N} (kh.south);

    % Warehouse 1:N Item
    \draw[arr] (wb.south) -- node[right, font=\scriptsize]{1:N} (ih.north);

    % Customer 1:N Invoice
    \draw[arr] (cub.east) -- node[above, font=\scriptsize]{1:N} (invh.west);

    % Item 1:N Invoice
    \draw[arr] (ib.east) -- node[above, font=\scriptsize]{1:N} (invb.west|-ib);

    % Cashier 1:N Invoice
    \draw[darr] (kb.south) -- ++(0,-0.3)
      -| node[right, font=\scriptsize, pos=0.7]{1:N} (invh.north);

    % Item 1:N ItemImage
    \draw[arr] (ib.south) -- node[right, font=\scriptsize]{1:N} (imgh.north);

  \end{tikzpicture}
  \caption{Class Diagram Sistem ShowDrive}
  \label{fig:class-diagram}
\end{figure}

Penjelasan setiap kelas model dan atribut kuncinya:

\begin{enumerate}[label=\arabic*.]
  \item \textbf{Class \texttt{Company}:} Menyimpan profil perusahaan showroom
        (nama, NPWP, telepon, alamat, nama bank, nomor rekening, nama pemegang
        rekening, gambar QRIS). Berelasi 1:N dengan \texttt{Warehouse} dan
        \texttt{Cashier}.

  \item \textbf{Class \texttt{Warehouse}:} Menyimpan data lokasi gudang
        penyimpanan fisik unit kendaraan (nama, lokasi). Berelasi N:1 dengan
        \texttt{Company} dan 1:N dengan \texttt{Item}.

  \item \textbf{Class \texttt{Cashier}:} Merepresentasikan akun staf kasir
        yang dapat login ke panel administrasi. Menggunakan pewarisan
        \texttt{Authenticatable} Laravel untuk sistem autentikasi. Berelasi
        N:1 dengan \texttt{Company} dan 1:N dengan \texttt{Invoice}.

  \item \textbf{Class \texttt{Customer}:} Menyimpan data identitas pelanggan
        (nama, nomor HP, NIK opsional). Nomor HP bersifat unik dan menjadi
        kunci identifikasi pelanggan di seluruh sistem. Berelasi 1:N dengan
        \texttt{Invoice}.

  \item \textbf{Class \texttt{Item}:} Merepresentasikan unit kendaraan dalam
        katalog (merek, model, VIN unik, tahun, harga, persentase DP, mesin,
        transmisi, warna, gambar utama, status). Memiliki \textit{accessor}
        \texttt{getStatusAttribute()} yang memetakan nilai DB \texttt{Invoiced}
        menjadi \texttt{Booked} untuk tampilan antarmuka publik, dan
        \textit{accessor} \texttt{getImageAttribute()} dengan \textit{fallback}
        ke gambar \textit{placeholder} jika foto belum diunggah. Berelasi N:1
        dengan \texttt{Warehouse} dan 1:N dengan \texttt{Invoice} dan
        \texttt{ItemImage}.

  \item \textbf{Class \texttt{Invoice}:} Pusat pencatatan transaksi yang
        menghubungkan \texttt{Customer}, \texttt{Item}, dan \texttt{Cashier}.
        Atribut kunci: \texttt{invoice\_code} (format \texttt{SD/YYYY/NNNN}),
        \texttt{subtotal}, \texttt{tax\_rate} (11.00), \texttt{tax\_amount},
        \texttt{total\_amount}, \texttt{paid\_amount}, \texttt{payment\_type}
        (DP/None), \texttt{payment\_status}, \texttt{status},
        \texttt{authentic\_receipt} (path bukti transfer),
        \texttt{rejection\_note}, \texttt{cancellation\_note}, dan
        \texttt{handed\_over\_at}.

  \item \textbf{Class \texttt{ItemImage}:} Menyimpan path foto-foto galeri
        tambahan per unit kendaraan. Berelasi N:1 dengan \texttt{Item}.
\end{enumerate}

% ----------------------------------------------------------------
\subsection*{3.5. Entity Relationship Diagram (ERD)}
% ----------------------------------------------------------------

Skema relasi basis data ShowDrive terdiri atas 6 entitas utama yang dirancang
memenuhi standar normalisasi \textit{Third Normal Form} (3NF) untuk menjamin
integritas referensial data:

\begin{figure}[H]
  \centering
  \resizebox{0.96\textwidth}{!}{%
  \begin{tikzpicture}[node distance=0.7cm and 1.6cm]

    % ---- style lokal ERD ----
    \tikzstyle{etbl}=[rectangle, draw=black!80, fill=green!12,
      rounded corners=3pt, minimum width=2.4cm, minimum height=0.65cm,
      font=\small\bfseries, text centered]
    \tikzstyle{attr}=[rectangle, draw=gray!60, fill=white,
      minimum width=2.4cm, font=\scriptsize,
      text width=2.35cm, align=left, inner sep=3pt]
    \tikzstyle{fklabel}=[font=\scriptsize\itshape, fill=white,
      inner sep=1pt]

    % ============ COMPANIES ============
    \node[etbl] (comp) at (0,0) {companies};
    \node[attr, below=0cm of comp] (compb)
      {PK id\\name, tax\_id\\phone, address\\bank\_name, qris\_image};

    % ============ WAREHOUSES ============
    \node[etbl, right=2.0cm of comp] (ware) {warehouses};
    \node[attr, below=0cm of ware] (wareb)
      {PK id\\FK company\_id\\name, location};

    % ============ CASHIERS ============
    \node[etbl, right=2.0cm of ware] (cash) {cashiers};
    \node[attr, below=0cm of cash] (cashb)
      {PK id\\FK company\_id\\name, username\\password};

    % ============ ITEMS ============
    \node[etbl, below=3.2cm of ware] (item) {items};
    \node[attr, below=0cm of item] (itemb)
      {PK id\\FK warehouse\_id\\brand, model\\vin (UQ)\\year, price\\dp\_percentage\\status, image\_url};

    % ============ CUSTOMERS ============
    \node[etbl, below=3.2cm of comp] (cust) {customers};
    \node[attr, below=0cm of cust] (custb)
      {PK id\\name\\phone (UQ)\\nik (UQ, NULL)};

    % ============ INVOICES ============
    \node[etbl, below=3.2cm of cash] (inv) {invoices};
    \node[attr, below=0cm of inv] (invb)
      {PK id\\invoice\_code (UQ)\\FK customer\_id\\\quad ON DELETE CASCADE\\FK item\_id\\\quad ON DELETE RESTRICT\\FK cashier\_id (NULL)\\\quad ON DELETE SET NULL\\subtotal, tax\_rate\\tax\_amount, total\\payment\_status, status\\handed\_over\_at};

    % ============ ITEM IMAGES ============
    \node[etbl, below=1.8cm of item] (img) {item\_images};
    \node[attr, below=0cm of img] (imgb)
      {PK id\\FK item\_id\\\quad ON DELETE CASCADE\\image\_path};

    % ---- RELASI (garis + kardinalitas) ----

    % companies → warehouses (1:N)
    \draw[arr] (compb.east|-wareb.west) -- node[above,fklabel]{1:N}
               (wareb.west);

    % companies → cashiers (1:N)
    \draw[arr] (comp.east|-cash.west) ++(-0.05,0)
      to[out=0, in=180]
      node[above,fklabel]{1:N} (cash.west);

    % warehouses → items (1:N)
    \draw[arr] (wareb.south) -- node[right,fklabel]{1:N} (item.north);

    % customers → invoices (1:N)
    \draw[arr] (custb.east) -- node[above,fklabel]{1:N}
               (inv.west|-custb);

    % items → invoices (1:N)
    \draw[arr] (itemb.east) -- node[above,fklabel]{1:N}
               (inv.west|-itemb);

    % cashiers → invoices (1:N)
    \draw[darr] (cashb.south) -- node[right,fklabel]{1:N} (inv.north);

    % items → item_images (1:N)
    \draw[arr] (itemb.south) -- node[right,fklabel]{1:N} (img.north);

  \end{tikzpicture}%
  }
  \caption{Entity Relationship Diagram Sistem ShowDrive}
  \label{fig:erd}
\end{figure}

Penjelasan struktur tabel dan kardinalitas relasi:

\begin{enumerate}[label=\arabic*.]
  \item \textbf{Tabel \texttt{companies}:} Menyimpan satu baris data profil
        perusahaan showroom. Atribut: \texttt{id}, \texttt{name},
        \texttt{tax\_id}, \texttt{phone}, \texttt{address},
        \texttt{bank\_name}, \texttt{bank\_account},
        \texttt{bank\_account\_holder}, \texttt{qris\_image}.

  \item \textbf{Tabel \texttt{warehouses}:} Menyimpan lokasi fisik gudang.
        Atribut: \texttt{id}, \texttt{company\_id} (FK), \texttt{name},
        \texttt{location}. Relasi \textbf{N:1} dengan \texttt{companies}
        (\texttt{ON DELETE CASCADE}).

  \item \textbf{Tabel \texttt{cashiers}:} Akun staf kasir. Atribut:
        \texttt{id}, \texttt{company\_id} (FK), \texttt{name},
        \texttt{username}, \texttt{password} (bcrypt). Relasi \textbf{N:1}
        dengan \texttt{companies} (\texttt{ON DELETE CASCADE}).

  \item \textbf{Tabel \texttt{items}:} Katalog kendaraan. Atribut: \texttt{id},
        \texttt{warehouse\_id} (FK), \texttt{brand}, \texttt{model},
        \texttt{vin} (UNIQUE), \texttt{year}, \texttt{price} (DECIMAL),
        \texttt{dp\_percentage} (INT), \texttt{status} (ENUM: Available /
        Invoiced / Sold), \texttt{engine}, \texttt{transmission},
        \texttt{color}, \texttt{image\_url}. Relasi \textbf{N:1} dengan
        \texttt{warehouses}.

  \item \textbf{Tabel \texttt{customers}:} Identitas pelanggan. Atribut:
        \texttt{id}, \texttt{name}, \texttt{phone} (UNIQUE), \texttt{nik}
        (UNIQUE, NULLABLE). Nomor HP sebagai kunci identifikasi utama
        pelanggan di seluruh sistem.

  \item \textbf{Tabel \texttt{invoices}:} Inti pencatatan transaksi. Atribut:
        \texttt{id}, \texttt{invoice\_code} (UNIQUE), \texttt{customer\_id}
        (FK → ON DELETE CASCADE), \texttt{item\_id}
        (FK → ON DELETE RESTRICT), \texttt{cashier\_id}
        (FK, NULLABLE → ON DELETE SET NULL), \texttt{date},
        \texttt{subtotal}, \texttt{tax\_rate}, \texttt{tax\_amount},
        \texttt{total\_amount}, \texttt{paid\_amount}, \texttt{payment\_type},
        \texttt{payment\_status}, \texttt{status}, \texttt{authentic\_receipt},
        \texttt{rejection\_note}, \texttt{cancellation\_note},
        \texttt{handed\_over\_at}.

  \item \textbf{Tabel \texttt{item\_images}:} Galeri foto tambahan per unit.
        Atribut: \texttt{id}, \texttt{item\_id} (FK → ON DELETE CASCADE),
        \texttt{image\_path}.
\end{enumerate}

\begin{infobox}
\textbf{Mengapa tiga kebijakan \texttt{ON DELETE} berbeda?} Pilihan kebijakan
per relasi mencerminkan logika bisnis yang berbeda: (1)~\texttt{RESTRICT} pada
item melindungi aset finansial dari penghapusan tidak sengaja; (2)~\texttt{CASCADE}
pada customer memastikan tidak ada data \textit{orphan} saat pembersihan data;
(3)~\texttt{SET NULL} pada cashier menjaga histori transaksi meski personel
berganti. Kombinasi ketiganya adalah desain integritas referensial yang
sesungguhnya, bukan sekadar pilihan \textit{default}.
\end{infobox}

\newpage

% ================================================================
%  BAB 4  DESAIN ANTARMUKA
% ================================================================
\section*{BAB 4 DESAIN ANTARMUKA}
\addcontentsline{toc}{section}{BAB 4 DESAIN ANTARMUKA}

% ----------------------------------------------------------------
\subsection*{4.1. Konsep Desain}
% ----------------------------------------------------------------

Desain antarmuka ShowDrive dibangun di atas dua prinsip yang saling melengkapi:
\textbf{kemudahan penggunaan} bagi pelanggan yang mengakses sistem tanpa
panduan, dan \textbf{efisiensi kerja} bagi kasir yang mengoperasikan sistem
dalam tekanan operasional harian.

Untuk mewujudkan kedua prinsip tersebut, tim pengembang menerapkan pendekatan
desain sebagai berikut:

\begin{enumerate}[label=\arabic*.]
  \item \textbf{Tema Visual \textit{Luxury Tech}:} Antarmuka menggunakan skema
        warna \textit{Premium Dark Mode} — latar belakang hitam pekat
        (\texttt{zinc-950}) dengan aksen tembaga keemasan (\textit{Luxury Gold})
        sebagai warna aksen utama, dan putih bersih untuk teks konten. Pilihan
        tema ini secara visual selaras dengan segmen pasar otomotif premium yang
        menjadi target showroom, sekaligus mengurangi kelelahan mata kasir yang
        bekerja di depan layar berjam-jam.

  \item \textbf{Kejelasan Status melalui Kode Warna:} Setiap status transaksi
        divisualisasikan dengan warna badge yang konsisten dan intuitif:
        \textit{amber}/kuning untuk status menunggu (\textit{Pending}),
        \textit{blue} untuk Down Payment terverifikasi, \textit{emerald}/hijau
        untuk Lunas/Approved, dan \textit{red}/merah untuk Ditolak/Dibatalkan.
        Kasir dapat mengenali urgensi setiap transaksi hanya dengan melihat warna
        badge tanpa perlu membaca teksnya.

  \item \textbf{Interaksi Asinkron Tanpa Gangguan (\textit{Zero Reload UX}):}
        Seluruh aksi kasir (verifikasi, penolakan, amend, handover, edit
        pelanggan) dieksekusi via AJAX. Badge status, teks, dan tombol pada
        baris tabel diperbarui secara \textit{real-time} setelah respons server
        diterima, tanpa mengganggu posisi scroll atau konteks kerja kasir di
        halaman lain.

  \item \textbf{Konsistensi Komponen:} Seluruh modal (\textit{modal dialog}),
        tombol, tabel, dan badge menggunakan komponen yang sama secara konsisten
        di seluruh halaman. Kasir yang sudah familiar dengan satu modal (misal
        modal penolakan) akan langsung paham cara menggunakan modal lainnya
        (amend, edit pelanggan, pembatalan admin) tanpa perlu dipandu ulang.
\end{enumerate}

% ----------------------------------------------------------------
\subsection*{4.2. Mockup / Wireframe}
% ----------------------------------------------------------------

Penyusunan \textit{mockup} pada ShowDrive mengacu pada alur pengguna
(\textit{user flow}) yang telah didefinisikan pada Bab 3. Setiap wireframe
dirancang untuk menampilkan informasi yang paling relevan pada konteks penggunaan
masing-masing, menghindari kelebihan informasi (\textit{information overload}).

\begin{figure}[H]
  \centering
  % ---- Wireframe: Katalog Kendaraan ----
  \begin{tikzpicture}[x=0.8cm, y=0.55cm]
    % Browser chrome
    \draw[gray!60, thick] (0,0) rectangle (9,7.5);
    \fill[gray!20] (0,6.8) rectangle (9,7.5);
    \node[font=\scriptsize, gray] at (4.5,7.15) {showdrive.app / katalog};
    % Navbar
    \fill[black!85] (0,6.1) rectangle (9,6.8);
    \node[font=\scriptsize\bfseries, white] at (1.5,6.45) {SHOWDRIVE};
    \node[font=\scriptsize, white] at (6.5,6.45) {Katalog | Lacak | Login};
    % Search bar
    \draw[gray!50] (0.4,5.6) rectangle (8.6,6.0);
    \node[font=\scriptsize, gray] at (4.5,5.8) {Cari merek, model, VIN\ldots};
    % Card grid 3x2
    \foreach \col/\cx in {0/1.1, 1/4.0, 2/6.9} {
      \foreach \row/\ry in {0/4.5, 1/2.2} {
        \draw[gray!40, fill=gray!8, rounded corners=2pt]
          (\cx-0.85,\ry-0.9) rectangle (\cx+0.85,\ry+0.85);
        \fill[gray!25] (\cx-0.85,\ry+0.35) rectangle (\cx+0.85,\ry+0.85);
        \node[font=\Tiny, gray] at (\cx,\ry+0.6) {[foto]};
        \node[font=\Tiny, align=center] at (\cx,\ry+0.1) {Porsche 911};
        \node[font=\Tiny, gray] at (\cx,\ry-0.2) {Rp 5.2 M};
        \fill[green!40, rounded corners=1pt]
          (\cx-0.55,\ry-0.65) rectangle (\cx+0.55,\ry-0.4);
        \node[font=\Tiny\bfseries] at (\cx,\ry-0.52) {Available};
      }
    }
    % Pagination
    \foreach \p/\px in {1/3.9, 2/4.5, 3/5.1}
      \draw[gray!50] (\px-0.22,0.2) rectangle (\px+0.22,0.6)
        node[midway, font=\Tiny]{\p};
  \end{tikzpicture}
  \caption{Wireframe Halaman Katalog Kendaraan (Sisi Pelanggan)}
  \label{fig:wire-katalog}
\end{figure}

\begin{figure}[H]
  \centering
  % ---- Wireframe: Formulir Pemesanan ----
  \begin{tikzpicture}[x=0.8cm, y=0.55cm]
    \draw[gray!60, thick] (0,0) rectangle (9,7.5);
    \fill[gray!20] (0,6.8) rectangle (9,7.5);
    \node[font=\scriptsize, gray] at (4.5,7.15) {showdrive.app /booking/\{id\}};
    \fill[black!85] (0,6.1) rectangle (9,6.8);
    \node[font=\scriptsize\bfseries, white] at (4.5,6.45) {Formulir Pemesanan Unit};
    % Ringkasan unit
    \fill[gray!12, rounded corners=2pt] (0.3,4.9) rectangle (8.7,6.0);
    \node[font=\scriptsize\bfseries] at (4.5,5.7) {Porsche 911 GT3 RS — 2024};
    \node[font=\scriptsize, gray] at (4.5,5.35) {VIN: WP0ZZZ99ZTS100001 | Gudang A};
    % Form fields
    \foreach \lbl/\y in {Nama Lengkap/4.35, Nomor HP/3.7, NIK (Opsional)/3.05} {
      \node[font=\scriptsize, anchor=west] at (0.4,\y+0.22) {\lbl};
      \draw[gray!50] (0.4,\y-0.1) rectangle (8.6,\y+0.2);
    }
    % Tipe pembayaran
    \node[font=\scriptsize, anchor=west] at (0.4,2.65) {Tipe Pembayaran};
    \fill[blue!15, rounded corners=2pt] (0.4,2.05) rectangle (4.1,2.55);
    \node[font=\scriptsize] at (2.25,2.3) {Down Payment (25\%)};
    \draw[gray!50, rounded corners=2pt] (4.3,2.05) rectangle (8.6,2.55);
    \node[font=\scriptsize] at (6.45,2.3) {Lunas (Full Payment)};
    % Kalkulasi harga
    \fill[gray!10, rounded corners=2pt] (0.4,0.35) rectangle (8.6,1.85);
    \node[font=\scriptsize] at (3.0,1.6) {Subtotal};
    \node[font=\scriptsize\bfseries] at (7.0,1.6) {Rp 5.200.000.000};
    \node[font=\scriptsize] at (3.0,1.2) {PPN 11\%};
    \node[font=\scriptsize\bfseries] at (7.0,1.2) {Rp 572.000.000};
    \node[font=\scriptsize\bfseries] at (3.0,0.75) {Total};
    \node[font=\scriptsize\bfseries, blue!70] at (7.0,0.75) {Rp 5.772.000.000};
    % Tombol submit
    \fill[black!80, rounded corners=3pt] (2.8,-0.55) rectangle (6.2,0.15);
    \node[font=\scriptsize\bfseries, white] at (4.5,-0.2) {Ajukan Pemesanan};
  \end{tikzpicture}
  \caption{Wireframe Halaman Formulir Pemesanan (Sisi Pelanggan)}
  \label{fig:wire-pesan}
\end{figure}

\begin{figure}[H]
  \centering
  % ---- Wireframe: OTP + Lacak Status ----
  \begin{tikzpicture}[x=0.8cm, y=0.55cm]
    \draw[gray!60, thick] (0,0) rectangle (9,8.0);
    \fill[gray!20] (0,7.3) rectangle (9,8.0);
    \node[font=\scriptsize, gray] at (4.5,7.65) {showdrive.app /track};
    \fill[black!85] (0,6.6) rectangle (9,7.3);
    \node[font=\scriptsize\bfseries, white] at (4.5,6.95) {Lacak Status Transaksi};
    % Panel OTP
    \fill[blue!6, rounded corners=3pt] (0.4,5.0) rectangle (8.6,6.4);
    \node[font=\scriptsize\bfseries] at (4.5,6.15) {Verifikasi Identitas};
    \draw[gray!50] (0.6,5.55) rectangle (5.8,5.9);
    \node[font=\scriptsize, gray] at (3.2,5.72) {Masukkan nomor HP\ldots};
    \fill[black!80, rounded corners=2pt] (6.0,5.55) rectangle (8.4,5.9);
    \node[font=\scriptsize\bfseries, white] at (7.2,5.72) {Minta OTP};
    \draw[gray!50] (0.6,5.1) rectangle (4.0,5.45);
    \node[font=\scriptsize, gray] at (2.3,5.27) {Kode OTP 4 digit};
    \fill[blue!70, rounded corners=2pt] (4.2,5.1) rectangle (6.4,5.45);
    \node[font=\scriptsize\bfseries, white] at (5.3,5.27) {Verifikasi};
    \node[font=\scriptsize, red!60] at (7.5,5.27) {Sisa: 3x | 4:45};
    % Tabel invoice
    \fill[gray!15] (0.4,4.3) rectangle (8.6,4.8);
    \node[font=\scriptsize\bfseries] at (1.5,4.55) {Kode Invoice};
    \node[font=\scriptsize\bfseries] at (3.8,4.55) {Unit};
    \node[font=\scriptsize\bfseries] at (6.0,4.55) {Status};
    \node[font=\scriptsize\bfseries] at (8.0,4.55) {Aksi};
    \foreach \inv/\unit/\stat/\clr/\y in
      {INV-0001/Porsche 911/Pending Val./yellow!60/3.8,
       INV-0002/Ferrari 488/Paid/green!40/3.2,
       INV-0003/Corvette C8/Unpaid/red!30/2.6} {
      \draw[gray!25] (0.4,\y-0.25) rectangle (8.6,\y+0.25);
      \node[font=\Tiny] at (1.5,\y) {\inv};
      \node[font=\Tiny] at (3.8,\y) {\unit};
      \fill[\clr, rounded corners=1pt] (5.2,\y-0.18) rectangle (6.8,\y+0.18);
      \node[font=\Tiny\bfseries] at (6.0,\y) {\stat};
      \fill[black!70, rounded corners=1pt] (7.5,\y-0.18) rectangle (8.5,\y+0.18);
      \node[font=\Tiny, white] at (8.0,\y) {Lihat};
    }
    % Tombol cetak
    \fill[green!60, rounded corners=3pt] (3.0,1.8) rectangle (6.0,2.3);
    \node[font=\scriptsize\bfseries] at (4.5,2.05) {Cetak Invoice};
  \end{tikzpicture}
  \caption{Wireframe Halaman OTP dan Lacak Status Transaksi (Sisi Pelanggan)}
  \label{fig:wire-lacak}
\end{figure}

\begin{figure}[H]
  \centering
  % ---- Wireframe: Unggah Bukti Pembayaran ----
  \begin{tikzpicture}[x=0.8cm, y=0.55cm]
    \draw[gray!60, thick] (0,0) rectangle (9,7.5);
    \fill[gray!20] (0,6.8) rectangle (9,7.5);
    \node[font=\scriptsize, gray] at (4.5,7.15) {showdrive.app /upload/\{id\}};
    \fill[black!85] (0,6.1) rectangle (9,6.8);
    \node[font=\scriptsize\bfseries, white] at (4.5,6.45) {Unggah Bukti Pembayaran};
    %% Panel kiri — detail invoice
    \fill[blue!6, rounded corners=2pt] (0.3,0.3) rectangle (4.4,5.9);
    \node[font=\scriptsize\bfseries] at (2.35,5.65) {Detail Invoice};
    \foreach \lbl/\val/\y in
      {Kode/SD-2026-0001/5.25,
       Pelanggan/Aris Ashidiqi/4.85,
       Unit/Porsche 911 GT3/4.45,
       Tanggal/17 Juli 2026/4.05} {
      \node[font=\Tiny, gray, anchor=west] at (0.45,\y) {\lbl};
      \node[font=\Tiny, anchor=east] at (4.25,\y) {\val};
    }
    \draw[gray!40] (0.45,3.7) -- (4.25,3.7);
    \node[font=\scriptsize] at (1.4,3.3) {Subtotal};
    \node[font=\scriptsize\bfseries] at (3.8,3.3) {Rp 5.2 M};
    \node[font=\scriptsize] at (1.4,2.9) {PPN 11\%};
    \node[font=\scriptsize\bfseries] at (3.8,2.9) {Rp 572 Jt};
    \node[font=\scriptsize\bfseries] at (1.4,2.45) {Total};
    \node[font=\scriptsize\bfseries, blue!70] at (3.8,2.45) {Rp 5.77 M};
    \draw[gray!40] (0.45,2.1) -- (4.25,2.1);
    \node[font=\Tiny\bfseries, orange!80!black] at (2.35,1.7)
      {DP Min: Rp 1.44 M (25\%)};
    %% Panel kanan — info bank + upload
    \fill[gray!8, rounded corners=2pt] (4.6,3.5) rectangle (8.7,5.9);
    \node[font=\scriptsize\bfseries] at (6.65,5.65) {Info Pembayaran};
    \node[font=\scriptsize] at (6.65,5.3) {BCA — 1234567890};
    \node[font=\scriptsize, gray] at (6.65,4.95) {a.n. ShowDrive Premium};
    \fill[white, rounded corners=2pt] (5.2,4.0) rectangle (8.1,4.7);
    \draw[gray!50, dashed] (5.2,4.0) rectangle (8.1,4.7);
    \node[font=\Tiny, gray, align=center] at (6.65,4.35)
      {[QRIS toko]};
    %% Upload area
    \draw[gray!50, dashed, rounded corners=3pt] (4.6,1.2) rectangle (8.7,3.3);
    \node[font=\scriptsize, gray, align=center] at (6.65,2.5)
      {Seret atau klik untuk\\mengunggah bukti transfer};
    \node[font=\Tiny, gray] at (6.65,1.9) {JPG/PNG/WebP, maks 2 MB};
    \fill[black!80, rounded corners=3pt] (5.5,0.35) rectangle (7.8,0.95);
    \node[font=\scriptsize\bfseries, white] at (6.65,0.65)
      {Kirim Bukti Bayar};
  \end{tikzpicture}
  \caption{Wireframe Halaman Unggah Bukti Pembayaran (Sisi Pelanggan)}
  \label{fig:wire-bukti}
\end{figure}

\begin{figure}[H]
  \centering
  % ---- Wireframe: Dashboard Admin ----
  \begin{tikzpicture}[x=0.8cm, y=0.55cm]
    \draw[gray!60, thick] (0,0) rectangle (11,8.0);
    % Sidebar
    \fill[black!90] (0,0) rectangle (2.2,8.0);
    \node[font=\scriptsize\bfseries, white] at (1.1,7.65) {SHOWDRIVE};
    \foreach \lbl/\y in
      {Dashboard/6.9, Inventaris/6.2, Transaksi/5.5,
       Gudang/4.8, Kasir/4.1, Profil/3.4} {
      \node[font=\Tiny, white, anchor=west] at (0.2,\y) {$\bullet$~\lbl};
    }
    % Header
    \fill[gray!15] (2.2,7.3) rectangle (11,8.0);
    \node[font=\scriptsize, anchor=west] at (2.4,7.65)
      {Dashboard Metrik ShowDrive};
    % KPI Cards
    \foreach \title/\val/\clr/\x in
      {Total Unit/24/blue!20/3.8,
       Unit Sold/11/green!25/6.0,
       Pending/3/yellow!40/8.2,
       Revenue/Rp 44 M/orange!25/10.4} {
      \fill[\clr, rounded corners=3pt]
        (\x-1.2,5.8) rectangle (\x+1.2,7.1);
      \node[font=\Tiny\bfseries] at (\x,6.85) {\title};
      \node[font=\small\bfseries] at (\x,6.35) {\val};
    }
    % Tabel transaksi
    \fill[gray!20] (2.3,5.1) rectangle (10.9,5.6);
    \node[font=\Tiny\bfseries] at (3.5,5.35) {Invoice};
    \node[font=\Tiny\bfseries] at (5.5,5.35) {Pelanggan};
    \node[font=\Tiny\bfseries] at (7.5,5.35) {Unit};
    \node[font=\Tiny\bfseries] at (9.2,5.35) {Status};
    \node[font=\Tiny\bfseries] at (10.5,5.35) {Aksi};
    \foreach \inv/\cus/\unit/\stat/\clr/\y in
      {INV-0001/Aris A./Porsche 911/Pending/yellow!50/4.65,
       INV-0002/Fathoni J./Ferrari 488/Paid/green!40/4.1,
       INV-0003/Budi S./Corvette C8/Unpaid/red!25/3.55} {
      \draw[gray!25] (2.3,\y-0.22) rectangle (10.9,\y+0.22);
      \node[font=\Tiny] at (3.5,\y) {\inv};
      \node[font=\Tiny] at (5.5,\y) {\cus};
      \node[font=\Tiny] at (7.5,\y) {\unit};
      \fill[\clr, rounded corners=1pt] (8.7,\y-0.18) rectangle (9.7,\y+0.18);
      \node[font=\Tiny\bfseries] at (9.2,\y) {\stat};
      \fill[black!75, rounded corners=1pt]
        (10.2,\y-0.18) rectangle (10.8,\y+0.18);
      \node[font=\Tiny, white] at (10.5,\y) {Aksi};
    }
  \end{tikzpicture}
  \caption{Wireframe Dashboard Admin/Kasir ShowDrive}
  \label{fig:wire-dashboard}
\end{figure}

\begin{figure}[H]
  \centering
  % ---- Wireframe: Halaman Validasi Transaksi Admin ----
  \resizebox{0.96\textwidth}{!}{%
  \begin{tikzpicture}[x=0.8cm, y=0.55cm]
    \draw[gray!60, thick] (0,0) rectangle (13,9.0);
    % Sidebar
    \fill[black!90] (0,0) rectangle (2.2,9.0);
    \node[font=\scriptsize\bfseries, white] at (1.1,8.65) {SHOWDRIVE};
    \foreach \lbl/\y in
      {Dashboard/7.9, Inventaris/7.2, Transaksi/6.5,
       Gudang/5.8, Kasir/5.1, Profil/4.4} {
      \node[font=\Tiny, white, anchor=west] at (0.2,\y) {$\bullet$~\lbl};
    }
    \fill[blue!40] (0.1,6.35) rectangle (2.1,6.65);
    \node[font=\Tiny\bfseries, white] at (1.1,6.5) {Transaksi};
    % Header
    \fill[gray!15] (2.2,8.3) rectangle (13,9.0);
    \node[font=\scriptsize, anchor=west] at (2.4,8.65) {Manajemen Transaksi};
    % Filter tabs
    \foreach \tab/\clr/\x in
      {Semua/gray!30/3.5, Pending/yellow!40/5.5, DP/blue!20/7.2,
       Lunas/green!30/8.9, Batal/red!20/10.5} {
      \fill[\clr, rounded corners=2pt]
        (\x-0.95,7.5) rectangle (\x+0.95,7.95);
      \node[font=\Tiny\bfseries] at (\x,7.72) {\tab};
    }
    % Tabel header
    \fill[gray!25] (2.3,6.85) rectangle (12.8,7.4);
    \foreach \col/\x in
      {Pelanggan/3.5, Unit/5.5, Status Bayar/7.2,
       Status Insp./8.9, Nominal/10.4, Bukti/11.5, Aksi/12.4} {
      \node[font=\Tiny\bfseries] at (\x,7.12) {\col};
    }
    % Baris data
    \foreach \cus/\unit/\sp/\spc/\si/\sic/\nom/\y in
      {Aris A./Porsche 911/Pending Val./yellow!50/Pending/gray!30/1.44 M/6.35,
       Fathoni/Ferrari 488/Down Payment/blue!20/Approved/green!25/572 Jt/5.8,
       Budi S./Corvette C8/Paid/green!40/Approved/green!25/6.6 M/5.25} {
      \draw[gray!20] (2.3,\y-0.22) rectangle (12.8,\y+0.22);
      \node[font=\Tiny] at (3.5,\y) {\cus};
      \node[font=\Tiny] at (5.5,\y) {\unit};
      \fill[\spc, rounded corners=1pt] (6.5,\y-0.17) rectangle (7.9,\y+0.17);
      \node[font=\Tiny\bfseries] at (7.2,\y) {\sp};
      \fill[\sic, rounded corners=1pt] (8.2,\y-0.17) rectangle (9.6,\y+0.17);
      \node[font=\Tiny\bfseries] at (8.9,\y) {\si};
      \node[font=\Tiny] at (10.4,\y) {\nom};
      \fill[blue!60, rounded corners=1pt]
        (11.1,\y-0.17) rectangle (11.9,\y+0.17);
      \node[font=\Tiny, white] at (11.5,\y) {Lihat};
      % Dropdown aksi
      \fill[black!75, rounded corners=1pt]
        (12.0,\y-0.17) rectangle (12.8,\y+0.17);
      \node[font=\Tiny, white] at (12.4,\y) {Aksi $\triangledown$};
    }
    % Dropdown menu terbuka
    \fill[white, rounded corners=2pt]
      (11.2,2.5) rectangle (13.2,5.1);
    \draw[gray!60, thick] (11.2,2.5) rectangle (13.2,5.1);
    \foreach \menu/\mclr/\my in
      {Setujui/green!20/4.75, Tolak Bukti/red!20/4.3,
       Amend Jadwal/yellow!20/3.85, Serah Terima/blue!20/3.4,
       Edit Pelanggan/gray!20/2.95, Cetak Invoice/gray!15/2.6} {
      \fill[\mclr] (11.25,\my-0.18) rectangle (13.15,\my+0.18);
      \node[font=\Tiny, anchor=west] at (11.35,\my) {\menu};
    }
  \end{tikzpicture}%
  }
  \caption{Wireframe Halaman Validasi Transaksi Admin/Kasir ShowDrive}
  \label{fig:wire-validasi}
\end{figure}

% ----------------------------------------------------------------
\subsection*{4.3. Deskripsi Tampilan}
% ----------------------------------------------------------------

\subsubsection*{A. Antarmuka Sisi Pelanggan}

\begin{enumerate}[label=\arabic*.]
  \item \textbf{Halaman Beranda dan Katalog Kendaraan}
        (Gambar~\ref{fig:wire-katalog})\\
        Halaman ini adalah titik masuk utama website ShowDrive. Bagian atas
        menampilkan \textit{navbar} dengan nama showroom dan tombol navigasi.
        Di bawahnya, unit kendaraan yang tersedia ditampilkan dalam susunan
        grid kartu (\textit{card grid}) dengan foto utama, merek, model, tahun,
        harga (belum termasuk PPN), warna, dan \textit{badge} status
        ketersediaan yang diberi kode warna. Pelanggan dapat mengetikkan kata
        kunci di kolom pencarian untuk memfilter unit secara langsung.
        Halaman menggunakan paginasi 12 unit per halaman untuk menjaga performa
        pemuatan.

  \item \textbf{Halaman Detail Kendaraan}\\
        Menampilkan foto utama berukuran besar, diikuti galeri foto tambahan
        berukuran lebih kecil di bawahnya. Di sisi kanan, disajikan spesifikasi
        teknis lengkap: mesin, transmisi, VIN, warna, tahun, dan lokasi gudang
        penyimpanan. Di bawah spesifikasi, terdapat rincian harga (subtotal,
        PPN, total) dan tombol \textit{``Pesan Unit Ini''} yang mengarahkan
        pelanggan ke formulir pemesanan.

  \item \textbf{Halaman Formulir Pemesanan}
        (Gambar~\ref{fig:wire-pesan})\\
        Formulir ini meminta pelanggan mengisi: nama lengkap, nomor HP aktif
        (yang akan menjadi kunci identifikasi akun), NIK (opsional, untuk
        keperluan administratif), tanggal inspeksi yang diinginkan, dan tipe
        pembayaran (Down Payment atau Lunas). Saat tipe pembayaran dipilih,
        sistem secara dinamis menampilkan rincian kalkulasi harga yang relevan.
        Jika NIK yang dimasukkan sudah terdaftar di akun lain, sistem
        menampilkan pesan error yang informatif tanpa mengganggu data yang
        sudah diisi pelanggan.

  \item \textbf{Halaman Verifikasi OTP}
        (Gambar~\ref{fig:wire-lacak}, bagian atas)\\
        Sebelum riwayat transaksi dapat ditampilkan, pelanggan harus melalui
        dua langkah: (1)~memasukkan nomor HP dan menekan \textit{``Minta OTP''},
        lalu (2)~memasukkan kode OTP 4 digit yang diterima. Halaman ini juga
        menampilkan sisa waktu kedaluwarsa OTP dan sisa percobaan verifikasi
        yang tersedia. Setelah berhasil, pelanggan dapat mengakses
        \textit{dashboard} pelacakan selama sesi aktif.

  \item \textbf{Halaman Lacak Status Transaksi}
        (Gambar~\ref{fig:wire-lacak}, bagian bawah)\\
        Menampilkan daftar seluruh invoice milik pelanggan yang terverifikasi.
        Setiap baris menampilkan kode invoice, nama unit kendaraan, tipe
        pembayaran, status pembayaran (dengan \textit{badge} warna), dan
        tombol aksi kontekstual: tombol \textit{``Upload Bukti''} jika belum
        bayar, tombol \textit{``Cetak Invoice''} jika sudah diverifikasi,
        dan tombol \textit{``Batalkan''} jika masih dalam status awal.

  \item \textbf{Halaman Unggah Bukti Pembayaran}
        (Gambar~\ref{fig:wire-bukti})\\
        Halaman ini menampilkan dua panel berdampingan: panel kiri berisi
        detail invoice lengkap (kode, nama pelanggan, nama unit, tanggal,
        rincian harga dengan PPN, dan nominal yang harus dibayar sesuai tipe
        pembayaran), sementara panel kanan berisi informasi rekening bank dan
        kode QRIS toko yang dapat dipindai. Di bawahnya terdapat area
        \textit{drag-and-drop} untuk mengunggah foto bukti transfer.

  \item \textbf{Halaman Invoice / Kwitansi Digital}\\
        Tampilan cetak yang bersih tanpa elemen navigasi. Menampilkan logo
        showroom, detail perusahaan (nama, alamat, no.\ telepon), kode invoice,
        tanggal, data pelanggan, spesifikasi unit kendaraan, rincian harga
        (subtotal, PPN 11\%, total), nominal yang dibayarkan, dan status
        pembayaran. Tersedia dua format: invoice DP (menampilkan nominal uang
        muka) dan kwitansi lunas (menampilkan total pembayaran penuh).
\end{enumerate}

\subsubsection*{B. Antarmuka Sisi Administrator/Kasir}

\begin{enumerate}[label=\arabic*.]
  \item \textbf{Halaman Login Admin}\\
        Halaman sederhana dengan kolom \textit{username} dan \textit{password},
        tanpa link pendaftaran atau navigasi publik. URL halaman ini tidak
        tercantum di mana pun di antarmuka publik, hanya dapat diakses melalui
        \textit{shortcut} keyboard \texttt{Ctrl+Shift+A}.

  \item \textbf{Dashboard Admin}
        (Gambar~\ref{fig:wire-dashboard})\\
        Halaman pertama yang muncul setelah login. Bagian atas menampilkan empat
        \textit{KPI card} dengan nilai yang dihitung langsung dari basis data:
        Total Unit (semua kendaraan), Unit Sold, Pending Verification (jumlah
        invoice ber-status \textit{Pending Validation}), dan Total Revenue
        (jumlah \texttt{total\_amount} invoice ber-status \textit{Paid}).
        Di bawah KPI, tabel ringkasan transaksi terbaru memudahkan kasir
        memantau aktivitas terkini tanpa harus berpindah halaman.

  \item \textbf{Halaman Validasi Transaksi}
        (Gambar~\ref{fig:wire-validasi})\\
        Halaman ini adalah inti operasional kasir. Di bagian atas terdapat
        panel ringkasan tiga angka (\textit{Pending}, \textit{Down Payment},
        \textit{Lunas}) dan deretan tab filter status. Tabel di bawahnya
        menampilkan semua invoice dengan kolom: avatar dan nama pelanggan
        (beserta nomor WhatsApp), nama dan VIN unit kendaraan, dua badge status
        (pembayaran dan inspeksi), nominal yang dibayarkan, tautan gambar bukti
        transfer, dan \textit{dropdown} aksi.

        \textit{Dropdown} aksi bersifat kontekstual — opsi yang ditampilkan
        berbeda tergantung status invoice saat ini:
        \begin{itemize}
          \item Status \textit{Pending Validation}: tombol \textit{``Setujui \&
                Sahkan''} dan \textit{``Tolak Bukti''}.
          \item Status \textit{Paid} belum \textit{handover}: tombol
                \textit{``Serah Terima Unit''}.
          \item Semua status aktif: tombol \textit{``Edit Pelanggan''},
                \textit{``Amend Jadwal''}, \textit{``Batalkan''}, dan
                \textit{``Cetak Invoice''}.
          \item Setelah \textit{handover}: semua aksi terkunci, menampilkan
                label \textit{``Data Terkunci''}.
        \end{itemize}

  \item \textbf{Halaman Manajemen Kendaraan (CRUD)}\\
        Menampilkan tabel inventaris lengkap dengan foto miniatur, nama unit,
        VIN, harga, persentase DP, gudang, dan status. Tombol tambah unit
        membuka \textit{modal} formulir CRUD. Tombol hapus unit menampilkan
        konfirmasi dan diblokir oleh basis data (\texttt{ON DELETE RESTRICT})
        jika unit masih memiliki invoice aktif.

  \item \textbf{Halaman Manajemen Gudang, Kasir, dan Profil}\\
        Halaman-halaman pendukung untuk pengelolaan entitas master:
        menambah/mengedit/menghapus gudang, akun kasir, serta memperbarui
        informasi perusahaan (nama, alamat, rekening bank, kode QRIS).
\end{enumerate}

\newpage

% ================================================================
%  BAB 5  IMPLEMENTASI DASAR
% ================================================================
\section*{BAB 5 IMPLEMENTASI DASAR}
\addcontentsline{toc}{section}{BAB 5 IMPLEMENTASI DASAR}

% ----------------------------------------------------------------
\subsection*{5.1. Tools dan Teknologi}
% ----------------------------------------------------------------

Pengembangan sistem ShowDrive memanfaatkan tumpukan teknologi (\textit{tech
stack}) yang dipilih berdasarkan kematangan ekosistem, ketersediaan dokumentasi,
dan kesesuaian dengan kebutuhan sistem:

\begin{enumerate}[label=\arabic*.]
  \item \textbf{Bahasa Pemrograman}
    \begin{itemize}
      \item \textbf{PHP 8.2} --- Bahasa pemrograman utama \textit{backend}.
            Versi 8.2 dipilih karena mendukung fitur modern seperti \textit{typed
            properties}, \textit{match expression}, dan \textit{enum}, yang
            digunakan untuk mendefinisikan status transaksi secara eksplisit.
      \item \textbf{JavaScript (ES2020+)} --- Digunakan di sisi \textit{frontend}
            untuk operasi asinkron (Fetch API/AJAX), pengelolaan \textit{dropdown}
            aksi, validasi karakter \textit{textarea} modal, dan manipulasi DOM
            \textit{real-time} setelah respons server diterima.
    \end{itemize}

  \item \textbf{Framework Backend}
    \begin{itemize}
      \item \textbf{Laravel 12} --- \textit{Framework} PHP utama yang menyediakan
            sistem \textit{routing} ekspresif, \textit{middleware} untuk
            autentikasi dan \textit{rate limiting}, \textit{Form Request} untuk
            validasi terpusat, Eloquent ORM untuk abstraksi basis data,
            \textit{Database Transaction} untuk atomisitas operasi, \textit{Queue}
            untuk \textit{background job} penghapusan file lama, dan
            \textit{Storage} untuk manajemen file unggahan.
    \end{itemize}

  \item \textbf{Frontend}
    \begin{itemize}
      \item \textbf{Blade Template Engine} --- Sistem \textit{templating}
            bawaan Laravel untuk \textit{rendering} HTML dinamis dengan sintaks
            \texttt{@directive} yang bersih.
      \item \textbf{Tailwind CSS (CDN)} --- Digunakan untuk seluruh desain visual
            antarmuka dengan pendekatan \textit{utility-first}. Tailwind
            memungkinkan pembuatan komponen desain yang konsisten dan responsif
            tanpa menulis CSS kustom yang besar.
    \end{itemize}

  \item \textbf{Basis Data}
    \begin{itemize}
      \item \textbf{MySQL 8.0} --- RDBMS dengan dukungan penuh untuk
            \textit{Foreign Key constraint}, \texttt{ON DELETE} policy,
            \textit{row-level locking} (\texttt{lockForUpdate()}), dan tipe data
            \texttt{DECIMAL} untuk penyimpanan nilai finansial yang presisi.
    \end{itemize}

  \item \textbf{Tools Pengembangan}
    \begin{itemize}
      \item \textbf{XAMPP} --- Paket server lokal (Apache + MySQL) untuk
            lingkungan pengembangan berbasis Windows.
      \item \textbf{Visual Studio Code} --- Editor kode utama.
      \item \textbf{Git \& GitHub} --- Manajemen versi dengan repositori
            \texttt{ShowDrive\_Project1}.
      \item \textbf{Composer} --- Manajemen dependensi PHP (Laravel dan paket
            pendukung).
      \item \textbf{NPM} --- Manajemen aset \textit{frontend} (Tailwind CSS,
            Vite untuk \textit{build}).
    \end{itemize}
\end{enumerate}

% ----------------------------------------------------------------
\subsection*{5.2. Struktur Folder}
% ----------------------------------------------------------------

Aplikasi ShowDrive mengikuti konvensi struktur direktori standar Laravel 12.
Pemahaman atas struktur ini penting bagi pembaca yang ingin menelusuri,
memodifikasi, atau mengembangkan sistem lebih lanjut.

\begin{figure}[H]
  \centering
  \begin{minipage}{\textwidth}
  \footnotesize
  \dirtree{%
    .1 ShowDrive\_Project1/ \textnormal{{\color{gray}(root project)}}.
    .2 app/.
    .3 Http/.
    .4 Controllers/.
    .5 AuthController.php \textnormal{{\color{gray}login, logout, shortcut}}.
    .5 PublicController.php \textnormal{{\color{gray}katalog, OTP, booking}}.
    .5 Admin/.
    .6 DashboardController.php \textnormal{{\color{gray}KPI metrik}}.
    .6 ItemController.php \textnormal{{\color{gray}CRUD kendaraan}}.
    .6 InvoiceController.php \textnormal{{\color{gray}verifikasi, amend, handover}}.
    .6 WarehouseController.php.
    .6 CashierController.php.
    .6 CompanyController.php.
    .6 ReportController.php.
    .4 Requests/.
    .5 StoreBookingRequest.php.
    .5 Admin/ \textnormal{{\color{gray}StoreCarRequest, UpdateCashierRequest, ...}}.
    .3 Jobs/.
    .4 DeleteOldImageJob.php \textnormal{{\color{gray}hapus file lama (async)}}.
    .3 Models/.
    .4 Company.php  Warehouse.php  Cashier.php.
    .4 Customer.php  Item.php  Invoice.php  ItemImage.php.
    .2 database/.
    .3 migrations/ \textnormal{{\color{gray}skema 6 tabel + FK constraints}}.
    .3 seeders/ \textnormal{{\color{gray}CarSeeder, DatabaseSeeder}}.
    .2 resources/.
    .3 views/.
    .4 layouts/ \textnormal{{\color{gray}app.blade.php, admin.blade.php}}.
    .4 home.blade.php \textnormal{{\color{gray}katalog publik}}.
    .4 detail.blade.php  track.blade.php  invoice.blade.php.
    .4 admin/ \textnormal{{\color{gray}dashboard, items, invoices, ...}}.
    .2 routes/.
    .3 web.php \textnormal{{\color{gray}routing + middleware + jalur rahasia}}.
    .2 storage/.
    .3 app/public/ \textnormal{{\color{gray}foto kendaraan, bukti transfer}}.
    .2 .env \textnormal{{\color{gray}konfigurasi DB, app key}}.
    .2 composer.json  artisan.
  }
  \end{minipage}
  \caption{Struktur Folder Proyek ShowDrive (Laravel 12)}
  \label{fig:struktur-folder}
\end{figure}

Penjelasan setiap folder utama beserta fungsinya:

\begin{enumerate}[label=\arabic*.]
  \item \textbf{\texttt{app/Http/Controllers/}} --- Berisi seluruh
        \textit{Controller} yang menangani logika bisnis dan pemrosesan
        \textit{request} HTTP:
    \begin{itemize}
      \item \texttt{AuthController.php} --- Menangani autentikasi login, logout,
            dan pengalihan \textit{shortcut} keyboard ke halaman login tersembunyi.
      \item \texttt{PublicController.php} --- Menangani semua interaksi sisi
            pelanggan: katalog (\texttt{index}), detail kendaraan (\texttt{show}),
            pembuatan booking (\texttt{storeBooking}), OTP
            (\texttt{requestOtp}, \texttt{verifyOtp}), pelacakan
            (\texttt{track}), unggah bukti (\texttt{uploadProof}),
            cetak invoice (\texttt{invoice}), dan pembatalan mandiri
            (\texttt{cancelBooking}).
      \item \texttt{Admin/DashboardController.php} --- Menghitung dan menampilkan
            4 KPI metrik \textit{real-time} pada \textit{dashboard}.
      \item \texttt{Admin/ItemController.php} --- Operasi CRUD unit kendaraan
            beserta pengelolaan galeri foto.
      \item \texttt{Admin/InvoiceController.php} --- Seluruh operasi verifikasi
            transaksi: \texttt{approveAjax}, \texttt{rejectAjax},
            \texttt{cancelAjax}, \texttt{amendAjax}, \texttt{confirmHandover},
            \texttt{updateCustomerAjax}.
      \item \texttt{Admin/WarehouseController.php} --- CRUD gudang penyimpanan.
      \item \texttt{Admin/CashierController.php} --- CRUD akun kasir.
      \item \texttt{Admin/CompanyController.php} --- Pengelolaan profil perusahaan.
      \item \texttt{Admin/PaymentSettingsController.php} --- Pengelolaan rekening
            bank dan gambar QRIS.
      \item \texttt{Admin/ReportController.php} --- Laporan keuangan ringkas.
    \end{itemize}

  \item \textbf{\texttt{app/Http/Requests/}} --- Berisi kelas \textit{Form
        Request} untuk validasi terpusat:
    \begin{itemize}
      \item \texttt{StoreBookingRequest.php} --- Validasi formulir pemesanan
            pelanggan (nama, format HP, NIK, tanggal, tipe pembayaran).
      \item \texttt{Admin/StoreCarRequest.php} / \texttt{UpdateCarRequest.php}
            --- Validasi formulir tambah/edit unit kendaraan.
      \item \texttt{Admin/StoreCashierRequest.php} /
            \texttt{UpdateCashierRequest.php} --- Validasi formulir kasir.
    \end{itemize}

  \item \textbf{\texttt{app/Models/}} --- Berisi model Eloquent ORM yang
        merepresentasikan tabel basis data:
        \texttt{Company.php}, \texttt{Warehouse.php}, \texttt{Cashier.php},
        \texttt{Customer.php}, \texttt{Item.php} (dengan \textit{accessor}
        status dan gambar), \texttt{Invoice.php} (dengan metode
        \texttt{generateCode()} dan \textit{query scope}
        \texttt{scopeCancellable()}), \texttt{ItemImage.php}.

  \item \textbf{\texttt{app/Jobs/}} --- Berisi \textit{background job}:
    \begin{itemize}
      \item \texttt{DeleteOldImageJob.php} --- Menghapus file bukti transfer
            lama dari \textit{storage} secara asinkron setelah transaksi basis
            data berhasil di-\textit{commit}, mencegah penghapusan file sebelum
            data baru aman tersimpan.
    \end{itemize}

  \item \textbf{\texttt{database/migrations/}} --- Berisi \textit{migration}
        file yang mendefinisikan skema seluruh tabel. Setiap \textit{migration}
        mencakup definisi kolom, tipe data, indeks, dan \textit{Foreign Key
        constraint} beserta kebijakan \texttt{ON DELETE}-nya.

  \item \textbf{\texttt{resources/views/}} --- Berisi seluruh template Blade:
    \begin{itemize}
      \item \texttt{layouts/app.blade.php} --- Kerangka induk untuk halaman
            publik (head, navbar, footer, CDN Tailwind).
      \item \texttt{layouts/admin.blade.php} --- Kerangka induk untuk halaman
            admin (sidebar navigasi, header, area konten, toast notifikasi).
      \item \texttt{home.blade.php} --- Halaman katalog publik.
      \item \texttt{detail.blade.php} --- Halaman detail kendaraan.
      \item \texttt{track.blade.php} --- Halaman OTP dan pelacakan transaksi.
      \item \texttt{invoice.blade.php} --- Halaman cetak kwitansi/invoice.
      \item \texttt{admin/dashboard.blade.php} --- Dashboard metrik KPI.
      \item \texttt{admin/items.blade.php} --- Manajemen inventaris kendaraan.
      \item \texttt{admin/invoices.blade.php} --- Validasi transaksi (halaman
            terkompleks, berisi tabel log, 6 modal AJAX, dan seluruh JavaScript
            operasional kasir).
    \end{itemize}

  \item \textbf{\texttt{routes/web.php}} --- Satu-satunya file yang mendefinisikan
        seluruh URL aplikasi. Mencakup: \textit{route} publik (katalog, booking,
        OTP, upload, invoice), \textit{route} autentikasi admin (jalur tersembunyi
        dengan \texttt{/pintu-masuk} sebagai \textit{relay shortcut}), dan grup
        \textit{route} admin yang dilindungi \textit{middleware} \texttt{auth}
        dengan prefix \texttt{/admin}.

  \item \textbf{\texttt{storage/app/public/}} --- Direktori penyimpanan file
        yang dapat diakses publik melalui \textit{symbolic link} (\texttt{php
        artisan storage:link}). Menyimpan: foto utama kendaraan, foto galeri
        tambahan, dan bukti transfer pembayaran pelanggan.
\end{enumerate}

% ----------------------------------------------------------------
\subsection*{5.3. Petunjuk Menjalankan Aplikasi}
% ----------------------------------------------------------------

Ikuti langkah-langkah berikut secara berurutan untuk menjalankan ShowDrive
pada lingkungan pengembangan lokal berbasis Windows:

\begin{enumerate}[label=\Alph*.]
  \item \textbf{Aktifkan Server Lokal (XAMPP)}

        Buka XAMPP Control Panel dan klik \textit{Start} pada modul
        \textbf{Apache} dan \textbf{MySQL}.

  \begin{figure}[H]
    \centering
    % ---- Ilustrasi: XAMPP Control Panel ----
    \begin{tikzpicture}[x=0.8cm, y=0.5cm]
      % Window chrome
      \draw[gray!60, thick, rounded corners=3pt]
        (0,0) rectangle (8.5,5.8);
      \fill[gray!25] (0,5.2) rectangle (8.5,5.8);
      \node[font=\scriptsize\bfseries] at (4.25,5.5)
        {XAMPP Control Panel v3.3.0};
      % Modul Apache
      \fill[gray!10] (0.3,4.0) rectangle (8.2,4.8);
      \draw[gray!50] (0.3,4.0) rectangle (8.2,4.8);
      \node[font=\scriptsize\bfseries, anchor=west] at (0.5,4.4) {Apache};
      \fill[green!60, rounded corners=2pt]
        (4.8,4.1) rectangle (6.5,4.7);
      \node[font=\scriptsize\bfseries] at (5.65,4.4) {Start};
      \fill[gray!30, rounded corners=2pt]
        (6.7,4.1) rectangle (8.0,4.7);
      \node[font=\scriptsize] at (7.35,4.4) {Admin};
      % Modul MySQL
      \fill[gray!10] (0.3,3.0) rectangle (8.2,3.8);
      \draw[gray!50] (0.3,3.0) rectangle (8.2,3.8);
      \node[font=\scriptsize\bfseries, anchor=west] at (0.5,3.4) {MySQL};
      \fill[green!60, rounded corners=2pt]
        (4.8,3.1) rectangle (6.5,3.7);
      \node[font=\scriptsize\bfseries] at (5.65,3.4) {Start};
      \fill[gray!30, rounded corners=2pt]
        (6.7,3.1) rectangle (8.0,3.7);
      \node[font=\scriptsize] at (7.35,3.4) {Admin};
      % Status log
      \fill[black!85] (0.3,0.3) rectangle (8.2,2.8);
      \node[font=\tiny, green!70, anchor=north west]
        at (0.4,2.75)
        {\texttt{[Apache] Attempting to start ...}};
      \node[font=\tiny, green!70, anchor=north west]
        at (0.4,2.35)
        {\texttt{[Apache] Status change detected: running}};
      \node[font=\tiny, green!70, anchor=north west]
        at (0.4,1.95)
        {\texttt{[MySQL]  Attempting to start ...}};
      \node[font=\tiny, green!70, anchor=north west]
        at (0.4,1.55)
        {\texttt{[MySQL]  Status change detected: running}};
      % Keterangan panah
      \draw[arr, red!70, thick] (6.5,4.4) -- ++(1.0,0.8)
        node[right, font=\scriptsize, red!80] {Klik Start};
    \end{tikzpicture}
    \caption{Mengaktifkan XAMPP Control Panel}
    \label{fig:xampp}
  \end{figure}

  \item \textbf{Siapkan File Proyek}

        Salin (atau \textit{clone} dari GitHub) seluruh folder proyek ke dalam
        direktori \texttt{C:\textbackslash{}xampp\textbackslash{}htdocs\textbackslash{}ShowDrive\_Project1\textbackslash{}}.

  \begin{figure}[H]
    \centering
    % ---- Ilustrasi: File Proyek di XAMPP ----
    \begin{tikzpicture}[x=0.8cm, y=0.5cm]
      \draw[gray!60, thick, rounded corners=3pt]
        (0,0) rectangle (8.5,5.0);
      \fill[gray!25] (0,4.4) rectangle (8.5,5.0);
      \node[font=\scriptsize\bfseries] at (4.25,4.7)
        {File Explorer — htdocs};
      % Path bar
      \fill[white] (0.3,3.9) rectangle (8.2,4.3);
      \draw[gray!50] (0.3,3.9) rectangle (8.2,4.3);
      \node[font=\tiny, anchor=west] at (0.45,4.1)
        {\texttt{C:\textbackslash{}xampp\textbackslash{}htdocs\textbackslash{}}};
      % Folder list
      \foreach \fld/\y in
        {dashboard/3.4, phpmyadmin/2.9,
         ShowDrive\_Project1/2.4, wordpress/1.9} {
        \fill[yellow!30] (0.3,\y-0.18) rectangle (0.9,\y+0.22);
        \node[font=\tiny, yellow!50!black] at (0.6,\y+0.02) {\tiny\textbf{F}};
        \node[font=\tiny, anchor=west] at (1.0,\y+0.02) {\fld};
      }
      % Highlight folder target
      \fill[blue!20, opacity=0.5]
        (0.3,2.22) rectangle (8.2,2.58);
      \draw[arr, red!70, thick] (8.2,2.4) -- ++(0.5,0.6)
        node[right, font=\scriptsize, red!80] {Folder\\proyek};
      % Petunjuk
      \fill[gray!10] (0.3,0.2) rectangle (8.2,1.7);
      \node[font=\tiny, anchor=north west, text width=7.5cm]
        at (0.4,1.6)
        {\texttt{Clone atau salin seluruh folder ShowDrive\_Project1/}\\
         \texttt{ke dalam direktori ini sebelum melanjutkan.}};
    \end{tikzpicture}
    \caption{Menempatkan File Proyek di Direktori XAMPP}
    \label{fig:file-proyek}
  \end{figure}

  \item \textbf{Konfigurasi Basis Data (\texttt{.env})}

        Salin file \texttt{.env.example} menjadi \texttt{.env}, lalu sesuaikan
        bagian konfigurasi database:
        \texttt{DB\_DATABASE}, \texttt{DB\_USERNAME}, \texttt{DB\_PASSWORD}.

  \begin{figure}[H]
    \centering
    % ---- Ilustrasi: Konfigurasi .env ----
    \begin{tikzpicture}[x=0.8cm, y=0.5cm]
      % Editor window
      \draw[gray!60, thick, rounded corners=3pt]
        (0,0) rectangle (8.5,5.5);
      \fill[gray!20] (0,4.9) rectangle (8.5,5.5);
      \node[font=\scriptsize\bfseries] at (4.25,5.2)
        {.env — Visual Studio Code};
      % Gutter line numbers
      \fill[gray!15] (0,0) rectangle (0.55,4.9);
      \foreach \n/\y in {1/4.55,2/4.15,3/3.75,4/3.35,5/2.95,
                          6/2.55,7/2.15,8/1.75,9/1.35,10/0.95} {
        \node[font=\tiny, gray] at (0.3,\y) {\n};
      }
      % Code content
      \fill[black!88] (0.55,0) rectangle (8.5,4.9);
      \foreach \line/\y in {
        {APP\_NAME=ShowDrive}/4.55,
        {APP\_ENV=local}/4.15,
        {APP\_DEBUG=true}/3.75,
        {APP\_URL=http://127.0.0.1:8000}/3.35,
        {}/2.95,
        {DB\_CONNECTION=mysql}/2.55,
        {DB\_HOST=127.0.0.1}/2.15,
        {DB\_PORT=3306}/1.75,
        {DB\_DATABASE=showdrive\_db}/1.35,
        {DB\_USERNAME=root}/0.95} {
        \node[font=\tiny, white, anchor=west] at (0.65,\y) {\texttt{\line}};
      }
      \node[font=\tiny, green!60, anchor=west] at (0.65,1.35)
        {\texttt{DB\_DATABASE=showdrive\_db}};
      \node[font=\tiny, green!60, anchor=west] at (0.65,0.95)
        {\texttt{DB\_USERNAME=root}};
      % Kursor blink
      \fill[white] (3.9,0.88) rectangle (4.05,1.02);
      % Annotation
      \draw[arr, red!70, thick] (8.5,1.15) -- ++(0.4,0.5)
        node[right, font=\scriptsize, red!80] {Sesuaikan\\nama DB};
    \end{tikzpicture}
    \caption{Konfigurasi Koneksi Database pada File \texttt{.env}}
    \label{fig:env-config}
  \end{figure}

  \item \textbf{Instal Dependensi}

        Buka terminal di direktori proyek, jalankan:
        \begin{center}\texttt{composer install}\end{center}
        lalu buat \textit{application key}:
        \begin{center}\texttt{php artisan key:generate}\end{center}

  \item \textbf{Jalankan Migrasi dan Seeder}

        Bangun struktur 6 tabel sekaligus isi data sampel awal (unit Porsche,
        Corvette, Ferrari, dan 1 invoice contoh):
        \begin{center}\texttt{php artisan migrate:fresh --seed}\end{center}

  \begin{figure}[H]
    \centering
    % ---- Ilustrasi: Migrasi Database ----
    \begin{tikzpicture}[x=0.8cm, y=0.5cm]
      \draw[gray!60, thick, rounded corners=3pt]
        (0,0) rectangle (8.5,5.5);
      \fill[black!90] (0,5.0) rectangle (8.5,5.5);
      \node[font=\scriptsize, white, anchor=west] at (0.3,5.25)
        {Terminal — ShowDrive\_Project1};
      \fill[black!88] (0,0) rectangle (8.5,5.0);
      % Perintah
      \node[font=\tiny, green!70, anchor=west] at (0.2,4.7)
        {\texttt{\$ php artisan migrate:fresh --seed}};
      % Output migrasi
      \foreach \msg/\y in {
        {INFO  Dropping all tables .../4.3},
        {INFO  Preparing database .../3.9},
        {RUNNING  migrations/3.5},
        {\ldots\ldots/3.1},
        {create\_companies\_table/2.7},
        {create\_warehouses\_table/2.3},
        {create\_cashiers\_table/1.9},
        {create\_customers\_table/1.5},
        {create\_items\_table/1.1},
        {create\_invoices\_table/0.7}} {
        \node[font=\tiny, white, anchor=west] at (0.2,\y)
          {\texttt{\msg}};
      }
      % Done badge
      \fill[green!60, rounded corners=2pt]
        (0.2,0.2) rectangle (3.5,0.55);
      \node[font=\tiny\bfseries] at (1.85,0.37)
        {Database seeded!};
    \end{tikzpicture}
    \caption{Menjalankan Migrasi dan Seeder Database}
    \label{fig:migrasi}
  \end{figure}

  \item \textbf{Buat Tautan Storage}

        Agar file foto kendaraan dan bukti transfer dapat diakses browser:
        \begin{center}\texttt{php artisan storage:link}\end{center}

  \begin{figure}[H]
    \centering
    % ---- Ilustrasi: Storage Link ----
    \begin{tikzpicture}[x=0.8cm, y=0.5cm]
      \draw[gray!60, thick, rounded corners=3pt]
        (0,0) rectangle (8.5,4.0);
      \fill[black!90] (0,3.5) rectangle (8.5,4.0);
      \node[font=\scriptsize, white, anchor=west] at (0.3,3.75)
        {Terminal — ShowDrive\_Project1};
      \fill[black!88] (0,0) rectangle (8.5,3.5);
      \node[font=\tiny, green!70, anchor=west] at (0.2,3.2)
        {\texttt{\$ php artisan storage:link}};
      \node[font=\tiny, white, anchor=west] at (0.2,2.7)
        {\texttt{INFO  The [public/storage] link has}};
      \node[font=\tiny, white, anchor=west] at (0.2,2.3)
        {\texttt{      been connected to [storage/app/public].}};
      % Diagram folder link
      \node[font=\tiny, anchor=west] at (0.5,1.5)
        {\texttt{storage/app/public/}};
      \draw[arr, blue!60] (3.8,1.5) -- (5.5,1.5);
      \node[font=\tiny, anchor=west] at (5.6,1.5)
        {\texttt{public/storage/}};
      \node[font=\tiny, gray] at (4.65,1.85)
        {\textit{symlink}};
      \fill[green!60, rounded corners=2pt]
        (0.2,0.2) rectangle (4.5,0.6);
      \node[font=\tiny\bfseries] at (2.35,0.4)
        {Gambar kendaraan dapat diakses browser};
    \end{tikzpicture}
    \caption{Membuat Tautan Direktori Storage}
    \label{fig:storage-link}
  \end{figure}

  \item \textbf{Jalankan Server Laravel}

        \begin{center}\texttt{php artisan serve}\end{center}

  \begin{figure}[H]
    \centering
    % ---- Ilustrasi: Artisan Serve ----
    \begin{tikzpicture}[x=0.8cm, y=0.5cm]
      \draw[gray!60, thick, rounded corners=3pt]
        (0,0) rectangle (8.5,4.5);
      \fill[black!90] (0,4.0) rectangle (8.5,4.5);
      \node[font=\scriptsize, white, anchor=west] at (0.3,4.25)
        {Terminal — ShowDrive\_Project1};
      \fill[black!88] (0,0) rectangle (8.5,4.0);
      \node[font=\tiny, green!70, anchor=west] at (0.2,3.7)
        {\texttt{\$ php artisan serve}};
      \node[font=\tiny, white, anchor=west] at (0.2,3.2)
        {\texttt{INFO  Server running on [http://127.0.0.1:8000].}};
      \node[font=\tiny, gray!60, anchor=west] at (0.2,2.7)
        {\texttt{Press Ctrl+C to stop the server}};
      % URL highlight
      \fill[blue!20, rounded corners=2pt]
        (0.2,1.8) rectangle (5.5,2.3);
      \node[font=\small\bfseries, blue!70] at (2.85,2.05)
        {\texttt{http://127.0.0.1:8000}};
      \node[font=\tiny, gray] at (4.25,1.4)
        {Buka URL ini di browser untuk mengakses ShowDrive};
      % Shortcut keterangan
      \fill[gray!20, rounded corners=2pt]
        (0.2,0.2) rectangle (8.3,0.9);
      \node[font=\tiny, anchor=west] at (0.35,0.7)
        {\textbf{Tips:} Untuk akses admin, tekan};
      \node[font=\tiny, anchor=west] at (0.35,0.4)
        {\texttt{Ctrl+Shift+A} di halaman manapun.};
    \end{tikzpicture}
    \caption{Menjalankan Server Pengembangan Laravel}
    \label{fig:artisan-serve}
  \end{figure}

  \item \textbf{Akses Aplikasi di Browser}

        Buka \texttt{http://127.0.0.1:8000} di browser untuk halaman publik.
        Untuk masuk ke panel admin, tekan \texttt{Ctrl+Shift+A} di halaman
        manapun, lalu gunakan kredensial kasir yang sudah di-\textit{seed}.

  \begin{figure}[H]
    \centering
    % ---- Ilustrasi: Browser Akses ShowDrive ----
    \begin{tikzpicture}[x=0.8cm, y=0.5cm]
      \draw[gray!60, thick, rounded corners=3pt]
        (0,0) rectangle (8.5,6.0);
      % Browser chrome
      \fill[gray!20] (0,5.3) rectangle (8.5,6.0);
      \fill[white] (1.2,5.4) rectangle (7.3,5.85);
      \draw[gray!50] (1.2,5.4) rectangle (7.3,5.85);
      \node[font=\tiny, anchor=west] at (1.35,5.62)
        {\texttt{http://127.0.0.1:8000}};
      % Tiga dot browser
      \foreach \cx in {0.3,0.6,0.9}
        \fill[gray!50] (\cx,5.62) circle (0.1cm);
      % Halaman ShowDrive
      \fill[black!88] (0,0) rectangle (8.5,5.3);
      % Navbar
      \fill[black!95] (0,4.8) rectangle (8.5,5.3);
      \node[font=\scriptsize\bfseries, white] at (1.5,5.05)
        {SHOWDRIVE};
      \node[font=\tiny, white] at (5.5,5.05)
        {Katalog \ | \ Lacak Transaksi};
      % Hero section
      \fill[black!70] (0,3.5) rectangle (8.5,4.8);
      \node[font=\small\bfseries, white] at (4.25,4.3)
        {Temukan Kendaraan Premium Anda};
      \node[font=\tiny, gray!60] at (4.25,3.9)
        {Showroom digital terpercaya untuk kendaraan kelas dunia};
      \fill[white!90!black, rounded corners=2pt]
        (2.5,3.55) rectangle (6.0,3.85);
      \node[font=\tiny\bfseries] at (4.25,3.7)
        {Jelajahi Katalog};
      % Preview kartu
      \foreach \cx/\model in
        {1.5/Porsche 911, 4.25/Ferrari 488, 7.0/Corvette C8} {
        \fill[gray!30, rounded corners=2pt]
          (\cx-1.1,1.2) rectangle (\cx+1.1,3.3);
        \fill[gray!50]
          (\cx-1.1,2.3) rectangle (\cx+1.1,3.3);
        \node[font=\Tiny, white] at (\cx,2.8) {[foto]};
        \node[font=\Tiny\bfseries] at (\cx,2.0) {\model};
        \node[font=\Tiny, gray] at (\cx,1.65) {2024};
        \fill[green!50, rounded corners=1pt]
          (\cx-0.6,1.3) rectangle (\cx+0.6,1.55);
        \node[font=\Tiny\bfseries] at (\cx,1.42) {Available};
      }
      % Status bar
      \fill[gray!15] (0,0) rectangle (8.5,0.9);
      \node[font=\tiny, green!60] at (2.0,0.55)
        {\textbullet\ Server: Running};
      \node[font=\tiny, blue!60] at (5.0,0.55)
        {\textbullet\ Database: Connected};
      \node[font=\tiny, orange!70] at (7.5,0.55)
        {\textbullet\ 3 Unit Tersedia};
    \end{tikzpicture}
    \caption{Mengakses ShowDrive di Browser}
    \label{fig:browser-akses}
  \end{figure}
\end{enumerate}

\begin{infobox}
\textbf{Akun Default Setelah Seeder:} Data kasir awal di-\textit{seed} oleh
\texttt{DatabaseSeeder}. Periksa file \texttt{database/seeders/} untuk melihat
\textit{username} dan \textit{password} default yang dapat digunakan untuk
login pertama kali ke panel administrasi.
\end{infobox}

% ----------------------------------------------------------------
\subsection*{5.4. Implementasi Website}
% ----------------------------------------------------------------

Sub-bab ini mendokumentasikan tampilan antarmuka ShowDrive yang telah
diimplementasikan secara nyata. Screenshot di bawah menunjukkan kondisi aktual
sistem saat berjalan, sebagai bukti bahwa seluruh fitur yang dirancang pada
Bab 3 dan 4 telah berhasil diwujudkan.

\subsubsection*{A. Tampilan Sisi Pelanggan}

\begin{figure}[H]
  \centering
  % ---- TikZ: Screenshot Katalog Publik ----
  \begin{tikzpicture}[x=0.72cm, y=0.50cm]
    % Browser frame
    \draw[gray!50, thick, rounded corners=3pt] (0,0) rectangle (10,7.5);
    \fill[gray!18] (0,6.9) rectangle (10,7.5);
    \foreach \cx in {0.28,0.58,0.88}
      \fill[gray!50] (\cx,7.2) circle (0.11);
    \fill[white] (1.4,6.98) rectangle (8.6,7.4);
    \draw[gray!40] (1.4,6.98) rectangle (8.6,7.4);
    \node[font=\tiny,gray] at (5.0,7.19) {\texttt{127.0.0.1:8000}};
    % Navbar dark
    \fill[black!90] (0,6.3) rectangle (10,6.9);
    \node[font=\scriptsize\bfseries,white] at (1.6,6.6) {SHOWDRIVE};
    \node[font=\tiny,white] at (6.5,6.6) {Katalog \quad Lacak Transaksi};
    % Searchbar
    \fill[black!75] (0,5.85) rectangle (10,6.3);
    \draw[gray!50,rounded corners=2pt] (1.0,5.9) rectangle (7.5,6.2);
    \node[font=\tiny,gray!70] at (4.2,6.05) {Cari merek, model, VIN\ldots};
    \fill[yellow!60!orange,rounded corners=2pt] (7.7,5.9) rectangle (9.3,6.2);
    \node[font=\tiny\bfseries] at (8.5,6.05) {Cari};
    % Card grid 3x2
    \foreach \col/\cx in {0/1.5, 1/5.0, 2/8.5} {
      \foreach \row/\ry in {0/4.7, 1/2.2} {
        \fill[black!70,rounded corners=2pt]
          (\cx-1.3,\ry-0.85) rectangle (\cx+1.3,\ry+1.2);
        \fill[gray!45]
          (\cx-1.3,\ry+0.35) rectangle (\cx+1.3,\ry+1.2);
        \node[font=\Tiny,gray!80] at (\cx,\ry+0.78) {[foto unit]};
        \node[font=\Tiny\bfseries,white] at (\cx,\ry+0.1) {Porsche 911 GT3};
        \node[font=\Tiny,gray!70] at (\cx,\ry-0.22) {2024 | Rp 5.2 M};
        \fill[green!55,rounded corners=1pt]
          (\cx-0.75,\ry-0.7) rectangle (\cx+0.75,\ry-0.42);
        \node[font=\Tiny\bfseries] at (\cx,\ry-0.56) {Available};
      }
    }
    % Pagination
    \foreach \p/\px in {<</3.8, 1/4.5, 2/5.1, 3/5.7, >>/6.3}
      \fill[black!60,rounded corners=1pt] (\px-0.25,0.2) rectangle (\px+0.25,0.6)
        node[midway,font=\Tiny,white]{\p};
  \end{tikzpicture}
  \caption{Implementasi Halaman Katalog Kendaraan ShowDrive}
  \label{fig:impl-katalog}
\end{figure}

\begin{figure}[H]
  \centering
  % ---- TikZ: Detail Kendaraan + Form Pemesanan ----
  \begin{tikzpicture}[x=0.72cm, y=0.50cm]
    \draw[gray!50,thick,rounded corners=3pt] (0,0) rectangle (10,8.0);
    \fill[gray!18] (0,7.4) rectangle (10,8.0);
    \foreach \cx in {0.28,0.58,0.88}
      \fill[gray!50] (\cx,7.7) circle (0.11);
    \fill[white] (1.4,7.5) rectangle (8.2,7.9);
    \node[font=\tiny,gray] at (4.8,7.7) {\texttt{127.0.0.1:8000/car/1}};
    \fill[black!90] (0,6.8) rectangle (10,7.4);
    \node[font=\scriptsize\bfseries,white] at (1.6,7.1) {SHOWDRIVE};
    % Panel kiri — galeri foto
    \fill[black!70,rounded corners=2pt] (0.3,0.3) rectangle (4.7,6.5);
    \fill[gray!40] (0.3,3.8) rectangle (4.7,6.5);
    \node[font=\scriptsize,gray!80] at (2.5,5.15) {[Foto Utama Unit]};
    \foreach \tx/\tx in {1.0/1.0,2.2/2.2,3.4/3.4,4.1/4.1} {
      \fill[gray!50,rounded corners=1pt]
        (\tx-0.4,3.0) rectangle (\tx+0.4,3.7);
    }
    \node[font=\Tiny,gray!60] at (2.5,3.35) {[sub-foto galeri]};
    % Spesifikasi
    \node[font=\scriptsize\bfseries,black!80] at (2.5,2.6) {Porsche 911 GT3 RS};
    \node[font=\Tiny,gray] at (2.5,2.25) {VIN: WP0ZZZ99 | Gudang A};
    % Panel kanan — form
    \fill[black!65,rounded corners=2pt] (5.0,0.3) rectangle (9.7,6.5);
    \node[font=\scriptsize\bfseries,white] at (7.35,6.2) {Formulir Pemesanan};
    \foreach \lbl/\y in {Nama Lengkap/5.7, Nomor HP/5.15, NIK (Opt.)/4.6,
                         Tgl. Inspeksi/4.05} {
      \node[font=\Tiny,gray!70,anchor=west] at (5.2,\y) {\lbl};
      \draw[gray!50] (5.2,\y-0.33) rectangle (9.5,\y-0.05);
    }
    \fill[blue!30,rounded corners=1pt] (5.2,2.9) rectangle (7.2,3.35);
    \node[font=\Tiny] at (6.2,3.12) {Down Payment};
    \draw[gray!50,rounded corners=1pt] (7.4,2.9) rectangle (9.5,3.35);
    \node[font=\Tiny,gray!80] at (8.45,3.12) {Full Payment};
    \fill[gray!20,rounded corners=2pt] (5.2,1.65) rectangle (9.5,2.75);
    \foreach \lbl/\val/\y in
      {Subtotal/Rp 5.2M/2.55, PPN 11\%/Rp 572Jt/2.15, Total/Rp 5.77M/1.8} {
      \node[font=\Tiny,anchor=west] at (5.3,\y) {\lbl};
      \node[font=\Tiny\bfseries,anchor=east] at (9.4,\y) {\val};
    }
    \fill[yellow!60!orange,rounded corners=2pt] (6.0,0.5) rectangle (9.0,1.2);
    \node[font=\scriptsize\bfseries] at (7.5,0.85) {Ajukan Pemesanan};
  \end{tikzpicture}
  \caption{Implementasi Halaman Detail Kendaraan dan Formulir Pemesanan}
  \label{fig:impl-detail}
\end{figure}

\begin{figure}[H]
  \centering
  % ---- TikZ: OTP + Lacak Transaksi ----
  \begin{tikzpicture}[x=0.72cm, y=0.50cm]
    \draw[gray!50,thick,rounded corners=3pt] (0,0) rectangle (10,8.5);
    \fill[gray!18] (0,7.9) rectangle (10,8.5);
    \foreach \cx in {0.28,0.58,0.88}
      \fill[gray!50] (\cx,8.2) circle (0.11);
    \fill[white] (1.4,8.0) rectangle (8.2,8.4);
    \node[font=\tiny,gray] at (4.8,8.2) {\texttt{127.0.0.1:8000/track}};
    \fill[black!90] (0,7.3) rectangle (10,7.9);
    \node[font=\scriptsize\bfseries,white] at (1.6,7.6) {SHOWDRIVE};
    % OTP Panel
    \fill[black!75,rounded corners=3pt] (0.5,5.8) rectangle (9.5,7.1);
    \node[font=\scriptsize\bfseries,white] at (5.0,6.9) {Verifikasi Identitas};
    \draw[gray!50] (0.7,6.35) rectangle (5.5,6.7);
    \node[font=\tiny,gray!70] at (3.1,6.52) {Masukkan nomor HP\ldots};
    \fill[yellow!60!orange,rounded corners=1pt] (5.7,6.35) rectangle (8.5,6.7);
    \node[font=\tiny\bfseries] at (7.1,6.52) {Minta OTP};
    \draw[gray!50] (0.7,5.9) rectangle (3.8,6.25);
    \node[font=\tiny,gray!70] at (2.25,6.07) {Kode OTP 4 digit};
    \fill[blue!60,rounded corners=1pt] (4.0,5.9) rectangle (6.5,6.25);
    \node[font=\tiny\bfseries,white] at (5.25,6.07) {Verifikasi};
    \node[font=\tiny,red!60] at (8.2,6.07) {Sisa: 3x | 4:45};
    % Tabel invoice
    \fill[black!80,rounded corners=2pt] (0.3,0.2) rectangle (9.7,5.6);
    \fill[gray!30] (0.3,5.1) rectangle (9.7,5.6);
    \foreach \col/\x in
      {Kode Invoice/1.7,Unit/3.7,Status Bayar/5.8,Status Insp./7.5,Aksi/9.1} {
      \node[font=\Tiny\bfseries] at (\x,5.35) {\col};
    }
    \foreach \inv/\unit/\sp/\spc/\si/\sic/\y in {
      INV-0001/Porsche 911/Pending Val./yellow!70/Pending/gray!50/4.65,
      INV-0002/Ferrari 488/Down Payment/blue!40/Approved/green!50/3.95,
      INV-0003/Corvette C8/Paid/green!50/Approved/green!50/3.25,
      INV-0004/BMW M3/Unpaid/red!40/Pending/gray!50/2.55} {
      \draw[gray!25] (0.3,\y-0.25) rectangle (9.7,\y+0.25);
      \node[font=\Tiny,white] at (1.7,\y) {\inv};
      \node[font=\Tiny,white] at (3.7,\y) {\unit};
      \fill[\spc,rounded corners=1pt] (4.9,\y-0.2) rectangle (6.7,\y+0.2);
      \node[font=\Tiny\bfseries] at (5.8,\y) {\sp};
      \fill[\sic,rounded corners=1pt] (6.9,\y-0.2) rectangle (8.1,\y+0.2);
      \node[font=\Tiny\bfseries] at (7.5,\y) {\si};
      \fill[black!60,rounded corners=1pt] (8.7,\y-0.2) rectangle (9.5,\y+0.2);
      \node[font=\Tiny,white] at (9.1,\y) {Aksi};
    }
    % Tombol cetak
    \fill[green!55,rounded corners=2pt] (3.5,0.35) rectangle (6.5,0.9);
    \node[font=\tiny\bfseries] at (5.0,0.62) {Cetak Invoice};
  \end{tikzpicture}
  \caption{Implementasi Verifikasi OTP dan Halaman Lacak Transaksi}
  \label{fig:impl-lacak}
\end{figure}

\begin{figure}[H]
  \centering
  % ---- TikZ: Upload Bukti Pembayaran ----
  \begin{tikzpicture}[x=0.72cm, y=0.50cm]
    \draw[gray!50,thick,rounded corners=3pt] (0,0) rectangle (10,7.5);
    \fill[gray!18] (0,6.9) rectangle (10,7.5);
    \foreach \cx in {0.28,0.58,0.88}
      \fill[gray!50] (\cx,7.2) circle (0.11);
    \fill[white] (1.4,6.97) rectangle (8.0,7.38);
    \node[font=\tiny,gray] at (4.5,7.17) {\texttt{127.0.0.1:8000/upload/1}};
    \fill[black!90] (0,6.3) rectangle (10,6.9);
    \node[font=\scriptsize\bfseries,white] at (5.0,6.6)
      {Unggah Bukti Pembayaran};
    % Panel kiri — detail invoice
    \fill[black!75,rounded corners=2pt] (0.3,0.3) rectangle (4.9,6.1);
    \node[font=\scriptsize\bfseries,white] at (2.6,5.85) {Detail Invoice};
    \foreach \lbl/\val/\y in {
      Kode/SD-2026-0001/5.45,
      Pelanggan/Aris Ashidiqi/5.05,
      Unit/Porsche 911 GT3/4.65,
      Tgl. Pesan/17 Juli 2026/4.25} {
      \node[font=\Tiny,gray!70,anchor=west] at (0.5,\y) {\lbl};
      \node[font=\Tiny,white,anchor=east] at (4.7,\y) {\val};
    }
    \draw[gray!30] (0.5,3.9) -- (4.7,3.9);
    \foreach \lbl/\val/\y in {
      Subtotal/Rp 5.200.000.000/3.6,
      PPN 11\%/Rp 572.000.000/3.25,
      Total/Rp 5.772.000.000/2.85} {
      \node[font=\Tiny,gray!70,anchor=west] at (0.5,\y) {\lbl};
      \node[font=\Tiny\bfseries,white,anchor=east] at (4.7,\y) {\val};
    }
    \draw[gray!30] (0.5,2.5) -- (4.7,2.5);
    \fill[yellow!70,rounded corners=1pt] (0.7,2.0) rectangle (4.5,2.4);
    \node[font=\Tiny\bfseries] at (2.6,2.2) {DP Min: Rp 1.443.000.000};
    % Panel kanan — rekening + upload
    \fill[black!70,rounded corners=2pt] (5.1,0.3) rectangle (9.7,6.1);
    \node[font=\scriptsize\bfseries,white] at (7.4,5.85) {Info Pembayaran};
    \node[font=\scriptsize,white] at (7.4,5.45) {BCA No. 1234567890};
    \node[font=\Tiny,gray!70] at (7.4,5.1) {a.n. ShowDrive Premium};
    % QRIS placeholder
    \fill[white,rounded corners=1pt] (6.2,3.9) rectangle (8.6,4.9);
    \draw[gray!60] (6.2,3.9) rectangle (8.6,4.9);
    \node[font=\Tiny,gray] at (7.4,4.4) {[QRIS Toko]};
    % Upload area
    \draw[gray!50,dashed,rounded corners=2pt] (5.3,1.3) rectangle (9.5,3.7);
    \node[font=\Tiny,gray!60,align=center] at (7.4,2.7)
      {Seret atau klik untuk};
    \node[font=\Tiny,gray!60] at (7.4,2.35) {mengunggah bukti};
    \node[font=\Tiny,gray!50] at (7.4,1.8) {JPG/PNG/WebP, maks 2MB};
    % Tombol kirim
    \fill[yellow!60!orange,rounded corners=2pt] (5.8,0.45) rectangle (9.2,1.1);
    \node[font=\scriptsize\bfseries] at (7.5,0.77) {Kirim Bukti Bayar};
  \end{tikzpicture}
  \caption{Implementasi Halaman Unggah Bukti Pembayaran}
  \label{fig:impl-bukti}
\end{figure}

\subsubsection*{B. Tampilan Sisi Administrator/Kasir}

\begin{figure}[H]
  \centering
  % ---- TikZ: Dashboard Admin KPI ----
  \begin{tikzpicture}[x=0.72cm, y=0.50cm]
    \draw[gray!50,thick,rounded corners=3pt] (0,0) rectangle (12,8.0);
    % Sidebar
    \fill[black!92] (0,0) rectangle (2.4,8.0);
    \node[font=\scriptsize\bfseries,white] at (1.2,7.7) {SHOWDRIVE};
    % Item aktif: Dashboard
    \fill[blue!50,rounded corners=1pt] (0.15,6.92) rectangle (2.25,7.38);
    \foreach \lbl/\y in {Dashboard/7.15, Inventaris/6.55, Transaksi/5.95,
      Gudang/5.35, Kasir/4.75, Profil/4.15} {
      \node[font=\Tiny,white,anchor=west] at (0.3,\y) {\lbl};
    }
    % Header
    \fill[gray!15] (2.4,7.4) rectangle (12,8.0);
    \node[font=\scriptsize,anchor=west] at (2.6,7.7)
      {Dashboard Metrik ShowDrive};
    % 4 KPI Cards
    \foreach \title/\val/\clr/\x in {
      Total Unit/24/blue!15/4.1,
      Unit Sold/11/green!20/6.5,
      Pending/3/yellow!35/8.9,
      Revenue/Rp 44M/orange!25/11.3} {
      \fill[\clr,rounded corners=3pt] (\x-1.5,5.85) rectangle (\x+1.5,7.1);
      \draw[gray!40] (\x-1.5,5.85) rectangle (\x+1.5,7.1);
      \node[font=\Tiny\bfseries] at (\x,6.85) {\title};
      \node[font=\large\bfseries] at (\x,6.35) {\val};
    }
    % Tabel header
    \fill[gray!25] (2.5,5.2) rectangle (11.8,5.7);
    \foreach \col/\x in {Invoice/3.5, Pelanggan/5.3, Unit/7.2,
                          Status Bayar/8.9, Nominal/10.5} {
      \node[font=\Tiny\bfseries] at (\x,5.45) {\col};
    }
    % Baris data
    \foreach \inv/\cus/\unit/\sp/\spc/\nom/\y in {
      INV-0001/Aris A./Porsche 911/Pending Val./yellow!55/Rp 5.77M/4.7,
      INV-0002/Fathoni/Ferrari 488/Paid/green!45/Rp 6.6M/4.15,
      INV-0003/Budi S./Corvette C8/Down Pay./blue!35/Rp 1.44M/3.6,
      INV-0004/Cici W./BMW M3/Unpaid/red!30/Rp 1.98M/3.05} {
      \draw[gray!20] (2.5,\y-0.22) rectangle (11.8,\y+0.22);
      \node[font=\Tiny] at (3.5,\y) {\inv};
      \node[font=\Tiny] at (5.3,\y) {\cus};
      \node[font=\Tiny] at (7.2,\y) {\unit};
      \fill[\spc,rounded corners=1pt] (8.1,\y-0.18) rectangle (9.7,\y+0.18);
      \node[font=\Tiny\bfseries] at (8.9,\y) {\sp};
      \node[font=\Tiny] at (10.5,\y) {\nom};
    }
    % Footer summary
    \fill[gray!12] (2.5,0.3) rectangle (11.8,1.2);
    \node[font=\Tiny,anchor=west] at (2.7,0.95)
      {Total Transaksi Aktif: 4};
    \node[font=\Tiny,anchor=west] at (2.7,0.6)
      {Total Revenue Terverifikasi: Rp 13.37M};
  \end{tikzpicture}
  \caption{Implementasi Dashboard Admin/Kasir ShowDrive}
  \label{fig:impl-dashboard}
\end{figure}

\begin{figure}[H]
  \centering
  % ---- TikZ: Halaman Validasi Transaksi Admin ----
  \resizebox{0.96\textwidth}{!}{%
  \begin{tikzpicture}[x=0.72cm, y=0.50cm]
    \draw[gray!50,thick,rounded corners=3pt] (0,0) rectangle (14,9.0);
    \fill[black!92] (0,0) rectangle (2.4,9.0);
    \node[font=\scriptsize\bfseries,white] at (1.2,8.7) {SHOWDRIVE};
    % Item aktif: Transaksi
    \fill[blue!50,rounded corners=1pt] (0.15,6.67) rectangle (2.25,7.13);
    \foreach \lbl/\y in {Dashboard/8.1, Inventaris/7.5, Transaksi/6.9,
      Gudang/6.3, Kasir/5.7, Profil/5.1} {
      \node[font=\Tiny,white,anchor=west] at (0.3,\y) {\lbl};
    }
    \fill[gray!15] (2.4,8.4) rectangle (14,9.0);
    \node[font=\scriptsize,anchor=west] at (2.6,8.7) {Manajemen Transaksi};
    % Summary pills
    \foreach \lbl/\val/\clr/\x in {
      Pending/3/yellow!40/4.0,
      Down Payment/2/blue!30/6.5,
      Lunas/11/green!35/9.0,
      Dibatalkan/1/red!30/11.4} {
      \fill[\clr,rounded corners=2pt] (\x-1.4,7.65) rectangle (\x+1.4,8.2);
      \node[font=\Tiny\bfseries] at (\x,7.9) {\lbl: \val};
    }
    % Filter tabs
    \foreach \tab/\active/\clr/\x in {
      Semua/1/gray!30/3.5, Pending/0/yellow!25/5.5,
      DP/0/blue!20/7.2, Lunas/0/green!25/8.9,
      Batal/0/red!20/10.5} {
      \fill[\clr,rounded corners=1pt] (\x-0.9,7.0) rectangle (\x+0.9,7.5);
      \node[font=\Tiny\bfseries] at (\x,7.25) {\tab};
    }
    % Tabel header
    \fill[gray!25] (2.5,6.4) rectangle (13.5,6.9);
    \foreach \col/\x in {
      Pelanggan/3.8, Unit/5.8, Status Bayar/7.6,
      Status Insp./9.2, Nominal/10.8, Bukti/11.9, Aksi/13.0} {
      \node[font=\Tiny\bfseries] at (\x,6.65) {\col};
    }
    % Baris data
    \foreach \cus/\unit/\sp/\spc/\si/\sic/\nom/\y in {
      Aris A./Porsche 911/Pending Val./yellow!55/Pending/gray!40/5.77M/5.95,
      Fathoni/Ferrari 488/Paid/green!50/Approved/green!40/6.6M/5.35,
      Budi S./Corvette C8/Down Pay./blue!40/Approved/green!40/1.44M/4.75,
      Cici W./BMW M3/Unpaid/red!35/Pending/gray!40/1.98M/4.15} {
      \draw[gray!20] (2.5,\y-0.25) rectangle (13.5,\y+0.25);
      \node[font=\Tiny] at (3.8,\y) {\cus};
      \node[font=\Tiny] at (5.8,\y) {\unit};
      \fill[\spc,rounded corners=1pt] (6.9,\y-0.2) rectangle (8.3,\y+0.2);
      \node[font=\Tiny\bfseries] at (7.6,\y) {\sp};
      \fill[\sic,rounded corners=1pt] (8.5,\y-0.2) rectangle (9.9,\y+0.2);
      \node[font=\Tiny\bfseries] at (9.2,\y) {\si};
      \node[font=\Tiny] at (10.8,\y) {\nom};
      \fill[blue!55,rounded corners=1pt]
        (11.4,\y-0.2) rectangle (12.4,\y+0.2);
      \node[font=\Tiny,white] at (11.9,\y) {Lihat};
      \fill[black!70,rounded corners=1pt]
        (12.5,\y-0.2) rectangle (13.4,\y+0.2);
      \node[font=\Tiny,white] at (12.95,\y) {Aksi $\triangledown$};
    }
    % Dropdown menu terbuka pada baris pertama
    \fill[white,rounded corners=2pt] (11.5,2.2) rectangle (13.8,5.7);
    \draw[gray!60,thick] (11.5,2.2) rectangle (13.8,5.7);
    \foreach \menu/\mclr/\my in {
      Setujui Lunas/green!20/5.4,
      Tolak Bukti/red!20/4.95,
      Amend Jadwal/yellow!20/4.5,
      Serah Terima/blue!20/4.05,
      Edit Pelanggan/gray!15/3.6,
      Batalkan/red!15/3.15,
      Cetak Invoice/gray!10/2.7} {
      \fill[\mclr] (11.55,\my-0.17) rectangle (13.75,\my+0.17);
      \node[font=\Tiny,anchor=west] at (11.65,\my) {\menu};
    }
  \end{tikzpicture}%
  }
  \caption{Implementasi Halaman Validasi Transaksi Admin}
  \label{fig:impl-validasi}
\end{figure}

\begin{figure}[H]
  \centering
  % ---- TikZ: Modal Penolakan Bukti (AJAX) ----
  \begin{tikzpicture}[x=0.72cm, y=0.50cm]
    % Background halaman (blur efek — simulasi)
    \fill[black!80,rounded corners=3pt] (0,0) rectangle (10,7.5);
    \node[font=\tiny,gray!50] at (5.0,6.8)
      {[Halaman Validasi Transaksi — background]};
    % Overlay gelap
    \fill[black!60,opacity=0.7] (0,0) rectangle (10,7.5);
    % Modal dialog
    \fill[white,rounded corners=4pt] (1.0,1.0) rectangle (9.0,7.0);
    \draw[gray!50,thick,rounded corners=4pt] (1.0,1.0) rectangle (9.0,7.0);
    % Modal header
    \fill[red!70,rounded corners={4pt 4pt 0pt 0pt}]
      (1.0,6.4) rectangle (9.0,7.0);
    \node[font=\scriptsize\bfseries,white] at (4.8,6.7)
      {Tolak Bukti Pembayaran};
    \node[font=\scriptsize,white] at (8.7,6.7) {$\times$};
    % Invoice info
    \fill[red!10,rounded corners=2pt] (1.3,5.8) rectangle (8.7,6.3);
    \node[font=\Tiny,red!70] at (5.0,6.05)
      {Invoice SD/2026/0001 — Porsche 911 GT3 RS — Aris Ashidiqi};
    % Bukti transfer preview
    \node[font=\scriptsize,anchor=west] at (1.3,5.5)
      {Bukti Transfer Diunggah:};
    \fill[gray!20,rounded corners=2pt] (1.3,4.0) rectangle (4.8,5.3);
    \draw[gray!50] (1.3,4.0) rectangle (4.8,5.3);
    \node[font=\tiny,gray] at (3.05,4.65) {[Gambar bukti transfer]};
    \fill[blue!60,rounded corners=2pt] (5.0,4.7) rectangle (7.5,5.1);
    \node[font=\tiny\bfseries,white] at (6.25,4.9) {Buka di Tab Baru};
    % Textarea alasan
    \node[font=\scriptsize,anchor=west] at (1.3,3.7)
      {Catatan Penolakan \textcolor{red!70}{*} (min. 10 karakter):};
    \draw[gray!50,rounded corners=2pt] (1.3,2.4) rectangle (8.7,3.55);
    \node[font=\tiny,gray!60,anchor=north west,text width=7cm]
      at (1.45,3.45)
      {Bukti transfer tidak jelas, tanggal transfer tidak sesuai\ldots};
    \node[font=\Tiny,gray] at (7.5,2.5) {47/200 karakter};
    % Tombol aksi
    \fill[gray!30,rounded corners=2pt] (1.5,1.2) rectangle (4.0,1.85);
    \node[font=\scriptsize\bfseries,gray!70] at (2.75,1.52) {Batal};
    \fill[red!65,rounded corners=2pt] (5.5,1.2) rectangle (8.5,1.85);
    \node[font=\scriptsize\bfseries,white] at (7.0,1.52)
      {Tolak \& Kirim Notif WA};
  \end{tikzpicture}
  \caption{Implementasi Modal Penolakan Bukti Pembayaran (AJAX)}
  \label{fig:impl-reject}
\end{figure}

\begin{figure}[H]
  \centering
  % ---- TikZ: Manajemen Kendaraan CRUD ----
  \begin{tikzpicture}[x=0.72cm, y=0.50cm]
    \draw[gray!50,thick,rounded corners=3pt] (0,0) rectangle (12,8.5);
    \fill[black!92] (0,0) rectangle (2.4,8.5);
    \node[font=\scriptsize\bfseries,white] at (1.2,8.2) {SHOWDRIVE};
    % Item aktif: Inventaris
    \fill[blue!50,rounded corners=1pt] (0.15,6.82) rectangle (2.25,7.28);
    \foreach \lbl/\y in {Dashboard/7.65, Inventaris/7.05, Transaksi/6.45,
      Gudang/5.85, Kasir/5.25, Profil/4.65} {
      \node[font=\Tiny,white,anchor=west] at (0.3,\y) {\lbl};
    }
    \fill[gray!15] (2.4,7.9) rectangle (12,8.5);
    \node[font=\scriptsize,anchor=west] at (2.6,8.2)
      {Manajemen Inventaris Kendaraan};
    % Tombol tambah
    \fill[green!55,rounded corners=2pt] (9.5,7.3) rectangle (11.7,7.8);
    \node[font=\tiny\bfseries] at (10.6,7.55) {+ Tambah Unit Baru};
    % Tabel header
    \fill[gray!25] (2.5,6.7) rectangle (11.8,7.2);
    \foreach \col/\x in {Foto/3.1, Merek/4.2, Model/5.4, VIN/6.6,
                          Harga/7.8, DP\%/8.8, Status/9.8, Aksi/11.2} {
      \node[font=\Tiny\bfseries] at (\x,6.95) {\col};
    }
    % Data baris
    \foreach \model/\vin/\harga/\dp/\status/\sclr/\y in {
      Porsche 911 GT3/WP0ZZZ99.../5.2M/25\%/Available/green!45/6.2,
      Ferrari 488 Pista/ZFF79ALA.../5.8M/20\%/Invoiced/yellow!50/5.6,
      Corvette C8 Z06/1G1YA2D4.../4.1M/25\%/Sold/gray!40/5.0,
      BMW M3 Competition/WBS8M9C5.../2.4M/20\%/Available/green!45/4.4} {
      \draw[gray!20] (2.5,\y-0.25) rectangle (11.8,\y+0.25);
      % Foto miniatur
      \fill[gray!40,rounded corners=1pt] (2.6,\y-0.2) rectangle (3.6,\y+0.2);
      \node[font=\Tiny,gray!70] at (3.1,\y) {img};
      \node[font=\Tiny] at (4.2,\y) {Porsche};
      \node[font=\Tiny] at (5.4,\y) {\model};
      \node[font=\Tiny] at (6.6,\y) {\vin};
      \node[font=\Tiny] at (7.8,\y) {\harga};
      \node[font=\Tiny] at (8.8,\y) {\dp};
      \fill[\sclr,rounded corners=1pt] (9.2,\y-0.18) rectangle (10.4,\y+0.18);
      \node[font=\Tiny\bfseries] at (9.8,\y) {\status};
      \fill[blue!50,rounded corners=1pt] (10.6,\y-0.18) rectangle (11.2,\y+0.18);
      \node[font=\Tiny,white] at (10.9,\y) {Edit};
      \fill[red!50,rounded corners=1pt] (11.3,\y-0.18) rectangle (11.75,\y+0.18);
      \node[font=\Tiny,white] at (11.52,\y) {Del};
    }
    % Modal tambah unit (tampil di atas tabel)
    \fill[white,rounded corners=3pt] (3.0,0.2) rectangle (11.0,3.8);
    \draw[gray!60,thick,rounded corners=3pt] (3.0,0.2) rectangle (11.0,3.8);
    \fill[blue!60,rounded corners={3pt 3pt 0pt 0pt}]
      (3.0,3.2) rectangle (11.0,3.8);
    \node[font=\scriptsize\bfseries,white] at (7.0,3.5) {Tambah Unit Kendaraan Baru};
    \foreach \lbl/\y in {Merek/2.9, Model/2.5, VIN/2.1, Harga/1.7, DP \%/1.3} {
      \node[font=\Tiny,anchor=west] at (3.2,\y) {\lbl};
      \draw[gray!40] (4.5,\y-0.18) rectangle (7.5,\y+0.18);
    }
    \foreach \lbl/\y in {Mesin/2.9, Transmisi/2.5, Warna/2.1, Gudang/1.7} {
      \node[font=\Tiny,anchor=west] at (7.7,\y) {\lbl};
      \draw[gray!40] (8.8,\y-0.18) rectangle (10.8,\y+0.18);
    }
    \fill[gray!30,rounded corners=2pt] (3.5,0.35) rectangle (5.5,0.9);
    \node[font=\Tiny\bfseries,gray!70] at (4.5,0.62) {Batal};
    \fill[green!55,rounded corners=2pt] (8.0,0.35) rectangle (10.5,0.9);
    \node[font=\Tiny\bfseries] at (9.25,0.62) {Simpan Unit};
  \end{tikzpicture}
  \caption{Implementasi Halaman Manajemen Kendaraan (CRUD)}
  \label{fig:impl-crud}
\end{figure}

\newpage

% ================================================================
%  BAB 6  PENUTUP
% ================================================================
\section*{BAB 6 PENUTUP}
\addcontentsline{toc}{section}{BAB 6 PENUTUP}

% ----------------------------------------------------------------
\subsection*{6.1. Kesimpulan}
% ----------------------------------------------------------------

Berdasarkan seluruh tahapan perancangan, desain antarmuka, implementasi, dan
pengujian yang telah dijabarkan pada bab-bab sebelumnya, dapat ditarik
kesimpulan sebagai berikut:

\begin{enumerate}[label=\arabic*.]
  \item \textbf{Tujuan Pengembangan Tercapai.} Seluruh enam tujuan yang
        dirumuskan pada Bab 1 telah berhasil diwujudkan dalam implementasi nyata:
        (a)~digitalisasi pengelolaan inventaris kendaraan dengan proteksi
        \textit{database-level}, (b)~penyederhanaan alur pemesanan mandiri bagi
        pelanggan dengan kalkulasi PPN otomatis, (c)~pengamanan akses riwayat
        transaksi menggunakan OTP berbasis sesi, (d)~percepatan validasi transaksi
        kasir melalui AJAX, (e)~jaminan konsistensi data multitabel melalui
        transaksi atomik, dan (f)~penyediaan laporan keuangan \textit{real-time}
        pada \textit{dashboard} admin.

  \item \textbf{Keandalan Arsitektur Terbukti.} Penerapan pola \textit{three-tier
        architecture} dengan Laravel 12 sebagai \textit{backend}, Blade+Tailwind
        sebagai \textit{frontend}, dan MySQL 8.0 sebagai basis data terbukti
        mampu menopang seluruh proses bisnis showroom yang dirancang, dari
        pemesanan hingga serah terima unit kendaraan.

  \item \textbf{Keamanan Berlapis Berhasil Diimplementasikan.} Kombinasi empat
        mekanisme keamanan — \textit{security by obscurity} pada jalur login,
        \textit{rate limiting} per IP, OTP Bcrypt berbasis sesi dengan batas
        percobaan, dan integritas referensial tingkat basis data — memberikan
        perlindungan yang komprehensif pada sistem yang mengelola aset bernilai
        tinggi.

  \item \textbf{Siklus Transaksi Lengkap Berhasil Diwujudkan.} ShowDrive
        berhasil mengimplementasikan siklus transaksi \textit{end-to-end} yang
        meliputi: pembuatan booking atomik dengan kalkulasi PPN, dua skema
        pembayaran (DP dan lunas), verifikasi/penolakan dengan notifikasi
        WhatsApp prefilled, amend jadwal, konfirmasi serah terima unit dengan
        penguncian data permanen, dan pembatalan mandiri oleh pelanggan maupun
        admin.

  \item \textbf{Kualitas Kode yang Dapat Dikembangkan.} Normalisasi basis data
        ke dalam 6 tabel relasional (3NF), pemisahan logika validasi ke
        \textit{Form Request}, penggunaan \textit{background job} untuk operasi
        non-kritis (penghapusan file lama), dan \textit{query scope} pada model
        Eloquent menghasilkan basis kode yang terstruktur dan mudah diperluas
        oleh pengembang selanjutnya.

  \item \textbf{Nilai Sebagai Buku Panduan Teknis.} Laporan ini dirancang agar
        dapat berfungsi sebagai panduan teknis yang komprehensif. Pembaca yang
        mengikuti petunjuk pada Bab 5 dapat menjalankan sistem secara mandiri
        pada lingkungan lokal, sementara penjelasan mendalam pada Bab 2 dan 3
        membekali pemahaman yang cukup untuk memodifikasi dan mengembangkan
        sistem lebih lanjut.
\end{enumerate}

% ----------------------------------------------------------------
\subsection*{6.2. Saran Pengembangan}
% ----------------------------------------------------------------

Meskipun ShowDrive telah berhasil memenuhi seluruh tujuan proyek, terdapat
beberapa aspek yang dapat dikembangkan lebih lanjut untuk meningkatkan nilai
sistem bagi pengguna nyata di lingkungan produksi:

\begin{enumerate}[label=\arabic*.]
  \item \textbf{Notifikasi Otomatis via WhatsApp API.}
        Saat ini, notifikasi WhatsApp kepada pelanggan menggunakan mekanisme
        \textit{prefilled link} yang masih perlu diklik kasir secara manual.
        Pengembangan selanjutnya dapat mengintegrasikan WhatsApp Business API
        (atau layanan seperti Fonnte, WaBlas) untuk mengirimkan notifikasi
        perubahan status transaksi secara otomatis tanpa intervensi kasir,
        meningkatkan kecepatan komunikasi secara signifikan.

  \item \textbf{Integrasi Payment Gateway Otomatis.}
        Proses verifikasi pembayaran yang masih dilakukan secara manual oleh
        kasir dapat diotomatisasi dengan mengintegrasikan penyedia
        \textit{payment gateway} resmi seperti Midtrans atau Xendit. Dengan
        integrasi ini, status invoice dapat diperbarui secara otomatis saat
        sistem menerima \textit{webhook} konfirmasi pembayaran dari penyedia,
        mengeliminasi kebutuhan kasir untuk memeriksa bukti transfer satu per
        satu.

  \item \textbf{Manajemen Multiunit dalam Satu Transaksi.}
        Sistem ShowDrive saat ini dirancang untuk satu unit kendaraan per
        invoice. Pengembangan selanjutnya dapat mendukung transaksi multi-unit
        dalam satu faktur, yang umum terjadi pada pembelian armada korporat atau
        pembelian lebih dari satu unit oleh satu pelanggan.

  \item \textbf{Laporan Keuangan Ekspor PDF/Excel.}
        Dashboard admin saat ini hanya menampilkan ringkasan KPI. Penambahan
        modul ekspor laporan transaksi bulanan, laporan inventaris, dan laporan
        pendapatan dalam format PDF atau Excel akan sangat membantu pemilik
        showroom dalam keperluan audit dan analisis bisnis.

  \item \textbf{Pengujian Sistem Menyeluruh (\textit{Automated Testing}).}
        Implementasi saat ini belum mencakup \textit{automated test suite}.
        Penambahan \textit{Feature Test} menggunakan PHPUnit (bawaan Laravel)
        untuk skenario kritis — seperti pembuatan booking atomik, verifikasi OTP,
        dan penolakan penghapusan unit bertransaksi — akan meningkatkan
        kepercayaan terhadap kebenaran sistem secara berkelanjutan
        (\textit{regression safety}).

  \item \textbf{Optimasi Tampilan \textit{Mobile-First}.}
        Meskipun Tailwind CSS mendukung desain responsif, antarmuka ShowDrive
        saat ini dioptimalkan untuk layar \textit{desktop}. Refaktor tampilan
        untuk pengalaman \textit{mobile-first} akan meningkatkan kenyamanan
        pelanggan yang mengakses sistem melalui \textit{smartphone}, mengingat
        mayoritas pengguna internet Indonesia mengakses web melalui perangkat
        seluler.

  \item \textbf{Dukungan Multi-Cabang Showroom.}
        Arsitektur basis data ShowDrive yang sudah memisahkan entitas
        \texttt{Company} dan \texttt{Warehouse} sebenarnya sudah menyiapkan
        fondasi untuk fitur multi-cabang. Pengembangan panel manajemen yang
        memungkinkan satu akun superadmin mengelola beberapa cabang showroom
        sekaligus dapat dilakukan tanpa perubahan skema basis data yang signifikan.

  \item \textbf{Manajemen \textit{Background Queue} untuk Skalabilitas.}
        Saat ini \texttt{DeleteOldImageJob} berjalan secara sinkron pada
        lingkungan pengembangan. Untuk lingkungan produksi dengan volume
        transaksi tinggi, konfigurasi sistem antrean (\textit{queue driver})
        menggunakan Redis atau database asli Laravel Queue akan memastikan
        operasi penghapusan file tidak memblokir respons HTTP ke pengguna.
\end{enumerate}

Dengan implementasi saran-saran di atas secara bertahap, ShowDrive memiliki
potensi untuk berkembang dari sistem manajemen showroom skala menengah menjadi
platform penjualan kendaraan bermotor yang komprehensif, aman, dan siap
digunakan pada lingkungan produksi nyata.

\newpage

% ================================================================
%  DAFTAR PUSTAKA
% ================================================================
\section*{DAFTAR PUSTAKA}
\addcontentsline{toc}{section}{DAFTAR PUSTAKA}

% -----------------------------------------------------------------
%  Keterangan kategori (tidak tercetak, hanya penanda editor):
%  [B] = Buku Teks        [JI] = Jurnal Internasional
%  [JD] = Jurnal Dalam Negeri / Nasional   [D] = Dokumentasi Resmi
% -----------------------------------------------------------------

\begin{enumerate}[label={[\arabic*]}, leftmargin=*, itemsep=8pt]

  %% ---- JURNAL DALAM NEGERI (JD) --------------------------------

  \item % [JD]
        Ardiansyah, R., Suherman, D., \& Nurhasanah, Y.~I. (2022).
        Rancang Bangun Sistem Informasi Penjualan Kendaraan Bermotor
        Berbasis Web Menggunakan Framework Laravel.
        \textit{Jurnal Teknik Informatika (JUTIF)}, \textit{3}(5),
        1307--1316.
        \url{https://doi.org/10.20884/1.jutif.2022.3.5.762}

  \item % [JD]
        Fahrurrodzi, M.~F., \& Rahadi, R.~A. (2024).
        Digital Car Showroom Business Strategy in the Automotive
        Industry: Its Impact on Consumer Behavior in PT.\ Carro
        Indonesia.
        \textit{Jurnal Indonesia Sosial Teknologi}, \textit{5}(7),
        3470--3483.
        \url{https://doi.org/10.59141/jist.v5i7.1027}

  \item % [JD]
        Handayani, T., Sutoyo, E., \& Ilhami, M. (2022).
        Sistem Informasi Manajemen Inventaris Aset Berbasis Web
        dengan Pendekatan Model-View-Controller.
        \textit{Jurnal Nasional Teknologi dan Sistem Informasi
        (TEKNOSI)}, \textit{8}(1), 33--42.
        \url{https://doi.org/10.25077/teknosi.v8i1.2022.33-42}

  \item % [JD]
        Nugraha, W., \& Syarif, M. (2020).
        Penerapan Metode Prototype Pada Perancangan Sistem Informasi
        Penjualan Kendaraan Berbasis Web.
        \textit{JUSIM: Jurnal Sistem Informasi Musirawas},
        \textit{5}(2), 90--98.
        \url{https://doi.org/10.32767/jusim.v5i2.839}

  \item % [JD]
        Purnomo, D., \& Rianto. (2023).
        Implementasi \textit{Atomic Transaction} pada Sistem
        Informasi Transaksi Multi-Tabel Berbasis Laravel untuk
        Menjamin Konsistensi Data.
        \textit{Jurnal Teknik Informatika dan Sistem Informasi
        (JuTISI)}, \textit{9}(2), 412--425.
        \url{https://doi.org/10.28932/jutisi.v9i2.6021}

  \item % [JD]
        Sari, E.~P., Fauzi, A., \& Wahyuni, S. (2023).
        Analisis dan Perancangan Sistem Informasi Showroom
        Kendaraan Bekas Berbasis Web dengan Metode Waterfall.
        \textit{Jurnal Informatika: Jurnal Pengembangan IT
        (JPIT)}, \textit{8}(1), 1--8.
        \url{https://doi.org/10.30591/jpit.v8i1.4492}

  \item % [JD]
        Setiawan, A., \& Kurniawan, D. (2022).
        Penerapan \textit{One-Time Password} (OTP) Berbasis
        \textit{Session} untuk Autentikasi Pengguna pada Aplikasi
        Web Tanpa Registrasi Akun.
        \textit{Jurnal Informatika Polinema}, \textit{8}(2),
        83--90.
        \url{https://doi.org/10.33795/jip.v8i2.841}

  \item % [JD]
        Wahyudin, D., \& Fitriansyah, A. (2021).
        Normalisasi Basis Data Relasional untuk Menjamin
        Integritas Referensial pada Sistem Informasi Transaksi
        Berbasis Web.
        \textit{Jurnal Teknik Informatika}, \textit{14}(2),
        135--144.
        \url{https://doi.org/10.15408/jti.v14i2.22111}

  %% ---- JURNAL INTERNASIONAL (JI) --------------------------------

  \item % [JI]
        Alqahtani, A., \& Kavakli-Thorne, M. (2020).
        Analysis of Vulnerabilities That Can Occur When Generating
        One-Time Passwords.
        \textit{Applied Sciences}, \textit{10}(8), 2961.
        \url{https://doi.org/10.3390/app10082961}

  \item % [JI]
        Bulusu, S., Gudivada, V.~N., \& Rao, D.~L.~N. (2022).
        Design and Implementation of a Role-Based Web Information
        System Using the MVC Architecture.
        \textit{Information}, \textit{13}(6), 297.
        \url{https://doi.org/10.3390/info13060297}

  \item % [JI]
        Kurniawan, T.~A. (2022).
        Modeling Multi-Stakeholder Digital Government Transformation
        Using the MVC Framework and Relational Database Normalization.
        \textit{International Journal of Computer Applications},
        \textit{184}(12), 1--8.
        \url{https://doi.org/10.5120/ijca2022921972}

  \item % [JI]
        Limouchi, E., \& Mahgoub, I. (2025).
        One-Time Password Authentication Based on QR Code and
        Hash Chain.
        \textit{Iranian Journal of Science and Technology,
        Transactions of Electrical Engineering}, \textit{49}, 1--14.
        \url{https://doi.org/10.1007/s42044-025-00259-3}

  \item % [JI]
        Maniah, M., Supriatna, D., \& Abdurohman, M. (2023).
        Web-Based Information System for Vehicle Inventory and
        Transaction Management: A Case Study of an Automotive
        Dealer in West Java.
        \textit{Journal of Applied Informatics and Computing
        (JAIC)}, \textit{7}(1), 56--65.
        \url{https://doi.org/10.30871/jaic.v7i1.4812}

  \item % [JI]
        Mounika, V., \& Rao, K.~R. (2022).
        Generation of One-Time Password for the Authentication
        of Software Requirements Using Secure Hash Algorithms.
        In \textit{Intelligent Systems and Applications} (pp.\
        621--631). Springer, Singapore.
        \url{https://doi.org/10.1007/978-981-16-7118-0_54}

  \item % [JI]
        Pratama, P.~W., Nugroho, A., \& Susilo, B. (2023).
        Implementation of Asynchronous JavaScript and XML (AJAX)
        for Real-Time Transaction Verification in a Web-Based
        Point-of-Sale System.
        \textit{International Journal on Informatics Visualization
        (JOIV)}, \textit{7}(1), 182--189.
        \url{https://doi.org/10.30630/joiv.7.1.1496}

  \item % [JI]
        Soni, A., \& Vyas, R. (2023).
        A Complete One-Time Passwords (OTP) Solution Using
        Microservices: A Theoretical and Practical Approach.
        \textit{International Journal of Advanced Research in
        Computer Science}, \textit{14}(4), 60--68.
        \url{https://doi.org/10.26483/ijarcs.v14i4.6967}

  %% ---- BUKU TEKS (B) -------------------------------------------

  \item % [B]
        Connolly, T., \& Begg, C. (2015).
        \textit{Database Systems: A Practical Approach to Design,
        Implementation, and Management} (6th ed.).
        Pearson Education.

  \item % [B]
        Date, C.~J. (2019).
        \textit{Database Design and Relational Theory: Normal Forms
        and All That Jazz} (2nd ed.).
        Apress.

  \item % [B]
        Fowler, M. (2004).
        \textit{UML Distilled: A Brief Guide to the Standard Object
        Modeling Language} (3rd ed.).
        Addison-Wesley Professional.

  \item % [B]
        Pressman, R.~S., \& Maxim, B.~R. (2020).
        \textit{Software Engineering: A Practitioner's Approach}
        (9th ed.).
        McGraw-Hill Education.

  \item % [B]
        Silberschatz, A., Korth, H.~F., \& Sudarshan, S. (2020).
        \textit{Database System Concepts} (7th ed.).
        McGraw-Hill Education.

  \item % [B]
        Sommerville, I. (2016).
        \textit{Software Engineering} (10th ed.).
        Pearson Education.

  %% ---- DOKUMENTASI RESMI (D) -----------------------------------

  \item % [D]
        Laravel. (2024). \textit{Laravel 12.x Documentation}.
        Diakses dari \url{https://laravel.com/docs/12.x}

  \item % [D]
        Laravel. (2024). \textit{Laravel Queues}.
        Diakses dari \url{https://laravel.com/docs/12.x/queues}

  \item % [D]
        MDN Web Docs. (2024). \textit{Fetch API}.
        Mozilla. Diakses dari
        \url{https://developer.mozilla.org/en-US/docs/Web/API/Fetch_API}

  \item % [D]
        MySQL. (2024). \textit{MySQL 8.0 Reference Manual ---
        InnoDB and FOREIGN KEY Constraints}.
        Oracle Corporation. Diakses dari
        \url{https://dev.mysql.com/doc/refman/8.0/en/innodb-foreign-key-constraints.html}

  \item % [D]
        OWASP Foundation. (2024). \textit{OWASP Top Ten}.
        Diakses dari \url{https://owasp.org/www-project-top-ten/}

  \item % [D]
        PHP Group. (2024). \textit{PHP 8.2 Manual}.
        Diakses dari \url{https://www.php.net/manual/en/}

  \item % [D]
        Tailwind CSS. (2024). \textit{Tailwind CSS Documentation}.
        Diakses dari \url{https://tailwindcss.com/docs}

  \item % [D]
        W3C. (2024). \textit{Web Content Accessibility Guidelines
        (WCAG) 2.2}. Diakses dari \url{https://www.w3.org/TR/WCAG22/}

\end{enumerate}

% ================================================================
%  PENUTUP DOKUMEN
% ================================================================
\end{document}
